% !TeX program = lualatex
\documentclass[12pt,aspectratio=169]{beamer}
\usepackage[utf8]{inputenc}
\usepackage[T1]{fontenc}
\usepackage{amsmath,amsfonts,amssymb}
\usepackage{graphicx}
\usepackage{xcolor}
\usepackage{hyperref}
\usepackage{anyfontsize}
\usepackage{lmodern}
\usepackage{longtable}
\usepackage{array}
\usepackage{booktabs}

% ====== EXTREME FONT REDUCTION ======
\renewcommand{\tiny}{\fontsize{3.5}{4.5}\selectfont}
\renewcommand{\scriptsize}{\fontsize{4.5}{5.5}\selectfont}
\renewcommand{\footnotesize}{\fontsize{5.5}{6.5}\selectfont}
\renewcommand{\small}{\fontsize{6.5}{7.5}\selectfont}
\renewcommand{\normalsize}{\fontsize{7.5}{8.5}\selectfont}
\renewcommand{\large}{\fontsize{8.5}{9.5}\selectfont}
\renewcommand{\Large}{\fontsize{9.5}{10.5}\selectfont}
\renewcommand{\LARGE}{\fontsize{10.5}{12}\selectfont}
\renewcommand{\huge}{\fontsize{11.5}{13.5}\selectfont}
\renewcommand{\Huge}{\fontsize{13}{16}\selectfont}

% ====== COLORS ======
\definecolor{redcorrect}{RGB}{200,30,30}
\definecolor{bluecorrect}{RGB}{30,80,200}
\definecolor{greencheck}{RGB}{0,150,50}
\definecolor{lightgray}{RGB}{245,245,245}
\definecolor{darkgray}{RGB}{40,40,40}
\definecolor{softyellow}{RGB}{255,248,200}
\definecolor{deepblue}{RGB}{20,60,150}
\definecolor{darkred}{RGB}{150,20,20}
\definecolor{orangewarning}{RGB}{255,140,0}
\definecolor{correctgreen}{RGB}{0,120,40}
\definecolor{trapcolor}{RGB}{180,80,0}
\definecolor{purpleconcept}{RGB}{120,50,180}
\definecolor{tealhard}{RGB}{0,120,120}

\setbeamercolor{title}{fg=white,bg=redcorrect}
\setbeamercolor{frametitle}{fg=white,bg=redcorrect}
\setbeamercolor{block title}{fg=white,bg=bluecorrect}
\setbeamercolor{block body}{fg=darkgray,bg=lightgray}
\setbeamercolor{normal text}{fg=darkgray}
\usecolortheme{default}

% ====== CUSTOM COMMANDS ======
\newcommand{\correct}[1]{\textcolor{greencheck}{\textbf{\checkmark}} #1}
\newcommand{\keypoint}[1]{\textcolor{deepblue}{\textbf{#1}}}
\newcommand{\hardconcept}[1]{\textcolor{darkred}{\textbf{#1}}}
\newcommand{\examtraps}[1]{\textcolor{trapcolor}{\textbf{⚠ TRAP:}} #1}
\newcommand{\recallbox}[1]{\fcolorbox{deepblue}{blue!5}{\parbox{0.88\textwidth}{\centering \textbf{REMEMBER:} #1}}}
\newcommand{\stepbox}[1]{%
	\begin{center}
		\fcolorbox{darkgray}{white}{%
			\parbox{0.92\textwidth}{\vspace{0.03cm} #1 \vspace{0.03cm}}%
		}
	\end{center}
}
\newcommand{\answerbox}[1]{\fcolorbox{correctgreen}{green!8}{\parbox{0.88\textwidth}{\centering \textcolor{correctgreen}{\textbf{\checkmark\ CORRECT:}} #1}}}
\newcommand{\wronganswer}[1]{\textcolor{redcorrect}{\textbf{✘}} #1}
\newcommand{\trapbox}[1]{\fcolorbox{trapcolor}{orange!8}{\parbox{0.88\textwidth}{\centering \textcolor{trapcolor}{\textbf{⚠ EXAM TRAP:}} #1}}}
\newcommand{\keyrule}[1]{\textcolor{tealhard}{\textbf{🔑}} #1}

\begin{document}
	
	% ================================================================
	% SLIDE 1: TITLE
	% ================================================================
	\begin{frame}
		\centering
		\vspace{0.5cm}
		{\fontsize{18}{22}\selectfont \textbf{CAMS EXAM MASTERY}}\\[0.15cm]
		{\fontsize{14}{17}\selectfont \textbf{ALL HARD QUESTIONS \& EXAM TRAPS}}\\[0.2cm]
		{\fontsize{9}{11}\selectfont \textit{Complete Tutorial — Every Concept You Get Wrong}}\\[0.1cm]
		\rule{4cm}{0.5pt}\\[0.1cm]
		{\fontsize{7}{9}\selectfont Authentication · VASP/MiCA · MLAT · PPP · FATF 6-Step · FSI}\\[0.05cm]
		{\fontsize{7}{9}\selectfont EWRA · Tax Evasion · Board Oversight · Terrorist Financing · DNFBPs}\\[0.05cm]
		{\fontsize{7}{9}\selectfont Layering · Structuring · IOSCO · IMF · Open Source · Batch Screening}\\[0.05cm]
		{\fontsize{7}{9}\selectfont Data Collection · Data Mining · Integration · BSA · China AML · AMLA}\\[0.05cm]
		{\fontsize{7}{9}\selectfont Concentration Accounts · Trade Topology · PSP · Casino · E-Commerce}\\[0.05cm]
		{\fontsize{7}{9}\selectfont Data Lineage · Clustering · Digital Onboarding · Perpetual KYC}\\[0.05cm]
		{\fontsize{7}{9}\selectfont Deepfake Detection · Entity Resolution · OFAC · Shell Companies}\\[0.05cm]
		{\fontsize{7}{9}\selectfont DORA · RPA · AI Transition · National e-ID · G-20 · Basel AML Index}
		\vspace{0.3cm}
	\end{frame}
	
	% ================================================================
	% SLIDE 2: AUTHENTICATION TECHNOLOGIES — THE BIOMETRICS TRAP
	% ================================================================
	\begin{frame}
		\frametitle{\fontsize{11}{13}\selectfont \textbf{Authentication — The Biometrics Trap}}
		\begin{block}{\fontsize{7}{9}\selectfont \textbf{The Hard Question — CAMS Exam Favorite}}
			\scriptsize
			Which technologies enhance authentication and security in financial systems? (Select Three.)
		\end{block}
		\vspace{0.02cm}
		\answerbox{\large \textbf{1. Multi-factor authentication. 2. Physiological biometrics. 3. Behavioral biometrics.}}
		\vspace{0.02cm}
		\begin{block}{\fontsize{7}{9}\selectfont \textbf{Natural Explanation — Why You Got This Wrong}}
			\scriptsize
			\begin{itemize}
				\item \wronganswer{Password biometrics} — This is a \textbf{trap}. There is \textbf{no such thing} as "password biometrics." Passwords are \textbf{not biometrics}.
				\item \textbf{Multi-factor authentication (MFA)} = multiple factors (something you know, have, are).
				\item \textbf{Physiological biometrics} = physical traits (fingerprint, facial recognition, iris scan, voice).
				\item \textbf{Behavioral biometrics} = behavioral patterns (typing rhythm, mouse movements, gait).
				\item \textbf{Key principle:} Biometrics are \textbf{physical or behavioral}. Passwords are \textbf{knowledge-based}.
			\end{itemize}
		\end{block}
		\vspace{0.02cm}
		\recallbox{\scriptsize \textbf{Key Rule:} MFA + physiological + behavioral = correct. "Password biometrics" = trap.}
	\end{frame}
	
	% ================================================================
	% SLIDE 3: VASP/MICA — THE WHITE PAPER TRAP
	% ================================================================
	\begin{frame}
		\frametitle{\fontsize{11}{13}\selectfont \textbf{VASP/MiCA — The White Paper Trap}}
		\begin{block}{\fontsize{7}{9}\selectfont \textbf{The Scenario — CAMS Exam Favorite}}
			\scriptsize
			A startup in Germany plans to issue a new crypto token and market it throughout the EU.
		\end{block}
		\vspace{0.02cm}
		\begin{block}{\fontsize{7}{9}\selectfont \textbf{The Hard Question}}
			\scriptsize
			What step must it take under MiCA regulations?
		\end{block}
		\vspace{0.02cm}
		\answerbox{\large \textbf{Publish a crypto-asset white paper and notify the national authority.}}
		\vspace{0.02cm}
		\begin{block}{\fontsize{7}{9}\selectfont \textbf{Natural Explanation — Why You Got This Wrong}}
			\scriptsize
			\begin{itemize}
				\item \wronganswer{Seek regulatory approval in each EU country} — This is what you'd do \textbf{without} MiCA. But MiCA creates a \textbf{harmonized} market.
				\item \textbf{White paper requirement:} Must publish a detailed disclosure document explaining the token, risks, and rights.
				\item \textbf{Notification:} Notify the national competent authority in your \textbf{home country} (Germany).
				\item \textbf{Passporting:} Once approved, can operate across \textbf{all EU countries} without additional approvals.
				\item \textbf{Key principle:} MiCA = \textbf{one notification + passporting}. NOT each country.
			\end{itemize}
		\end{block}
		\vspace{0.02cm}
		\recallbox{\scriptsize \textbf{Key Rule:} MiCA = \textbf{white paper + one national notification}. Passporting across EU.}
	\end{frame}
	
	% ================================================================
	% SLIDE 4: MLAT — THE ADMISSIBILITY TRAP
	% ================================================================
	\begin{frame}
		\frametitle{\fontsize{11}{13}\selectfont \textbf{MLAT — The Admissibility Trap}}
		\begin{block}{\fontsize{7}{9}\selectfont \textbf{The Scenario — CAMS Exam Favorite}}
			\scriptsize
			Two countries, A and B, are establishing an MLAT. Country A has \textbf{data protection rules} preventing sharing of client identities. Country B bases investigations on \textbf{beneficial ownership disclosures}.
		\end{block}
		\vspace{0.02cm}
		\begin{block}{\fontsize{7}{9}\selectfont \textbf{The Hard Question}}
			\scriptsize
			What challenge is most likely to arise in implementing the MLAT?
		\end{block}
		\vspace{0.02cm}
		\answerbox{\large \textbf{Evidence exchanged under the MLAT will NOT be admissible in Country B's legal system.}}
		\vspace{0.02cm}
		\begin{block}{\fontsize{7}{9}\selectfont \textbf{Natural Explanation — Why You Got This Wrong}}
			\scriptsize
			\begin{itemize}
				\item Most candidates think: \wronganswer{Country A's data protection rules will prevent sharing} — This is \textbf{obvious}. The exam tests \textbf{deeper} understanding.
				\item \textbf{Admissibility} is the legal standard for evidence in court.
				\item Even if Country A \textbf{shares} the data, Country B's courts may \textbf{reject} it.
				\item Why? Country B's evidentiary standards may require \textbf{different handling}, \textbf{chain of custody}, or \textbf{authentication}.
				\item \textbf{Key principle:} MLATs must address both \textbf{sharing} AND \textbf{admissibility}. The exam tests admissibility.
			\end{itemize}
		\end{block}
		\vspace{0.02cm}
		\recallbox{\scriptsize \textbf{Key Rule:} MLAT challenge = \textbf{admissibility}, not just sharing.}
	\end{frame}
	
	% ================================================================
	% SLIDE 5: PPP — THE REAL-TIME INTELLIGENCE TRAP
	% ================================================================
	\begin{frame}
		\frametitle{\fontsize{11}{13}\selectfont \textbf{PPP — The Real-Time Intelligence Trap}}
		\begin{block}{\fontsize{7}{9}\selectfont \textbf{The Scenario — CAMS Exam Trap}}
			\scriptsize
			An FIU wants to use the PPP model for \textbf{real-time intelligence sharing} to disrupt a terrorist financing network.
		\end{block}
		\vspace{0.02cm}
		\begin{block}{\fontsize{7}{9}\selectfont \textbf{The Hard Question}}
			\scriptsize
			Which approach is most aligned with PPP goals?
		\end{block}
		\vspace{0.02cm}
		\answerbox{\large \textbf{Establishing a formal platform with direct two-way communication on specific threats.}}
		\vspace{0.02cm}
		\begin{block}{\fontsize{7}{9}\selectfont \textbf{Natural Explanation — Why You Got This Wrong}}
			\scriptsize
			\begin{itemize}
				\item Many candidates choose \wronganswer{anonymized, aggregated format after threat is resolved} because it \textbf{sounds safer} or more \textbf{privacy-protecting}.
				\item But \textbf{that's not a PPP}. That's \textbf{after-the-fact reporting}.
				\item \textbf{PPP is about DISRUPTION} — stopping crimes \textbf{before} they happen.
				\item To disrupt, you need \textbf{real-time two-way communication}.
				\item Anonymized reporting \textbf{protects privacy} but \textbf{doesn't enable disruption}.
				\item \textbf{Key principle:} PPPs are about \textbf{real-time} intelligence sharing to \textbf{disrupt} active threats.
			\end{itemize}
		\end{block}
		\vspace{0.02cm}
		\recallbox{\scriptsize \textbf{Key Rule:} PPP = \textbf{real-time two-way communication} for disruption.}
	\end{frame}
	
	% ================================================================
	% SLIDE 6: FATF 6-STEP — THE COMPONENT TRAP
	% ================================================================
	\begin{frame}
		\frametitle{\fontsize{11}{13}\selectfont \textbf{FATF 6-Step — The Component Trap}}
		\begin{block}{\fontsize{7}{9}\selectfont \textbf{The Hard Question — CAMS Exam Favorite}}
			\scriptsize
			Which are key components of FATF's six-step framework for national ML/TF risk assessments? (Select Three.)
		\end{block}
		\vspace{0.02cm}
		\answerbox{\large \textbf{1. Environmental scanning. 2. Horizon scanning. 3. Threat identification.}}
		\vspace{0.02cm}
		\begin{block}{\fontsize{7}{9}\selectfont \textbf{Natural Explanation — Why You Got This Wrong}}
			\scriptsize
			\begin{itemize}
				\item \wronganswer{Cross-border tax enforcement} — This \textbf{sounds plausible} because it's related to financial crime, but it's \textbf{not} a FATF risk assessment component.
				\item \wronganswer{Offshore asset repatriation} — Also sounds related, but \textbf{not} a component.
				\item \textbf{Environmental scanning:} Understanding the context (what's happening in the country? Economic, political, social).
				\item \textbf{Horizon scanning:} Identifying emerging risks (what's coming? New technologies, new methods).
				\item \textbf{Threat identification:} Identifying specific threats (who are the bad actors? What are their methods?).
				\item \textbf{Key principle:} FATF's framework is about \textbf{understanding risk}, not enforcement actions.
			\end{itemize}
		\end{block}
		\vspace{0.02cm}
		\recallbox{\scriptsize \textbf{Key Rule:} FATF 6-step = \textbf{environmental + horizon scanning + threat identification}.}
	\end{frame}
	
	% ================================================================
	% SLIDE 7: FINANCIAL SECRECY INDEX — THE FACTORS TRAP
	% ================================================================
	\begin{frame}
		\frametitle{\fontsize{11}{13}\selectfont \textbf{Financial Secrecy Index — The Factors Trap}}
		\begin{block}{\fontsize{7}{9}\selectfont \textbf{The Hard Question — CAMS Exam Favorite}}
			\scriptsize
			Which contribute to an elevated Financial Secrecy Index rating? (Select Three.)
		\end{block}
		\vspace{0.02cm}
		\answerbox{\large \textbf{1. Minimal BO reporting. 2. Extensive cross-border services to non-residents. 3. High secrecy-to-scale ratio.}}
		\vspace{0.02cm}
		\begin{block}{\fontsize{7}{9}\selectfont \textbf{Natural Explanation — Why You Got This Wrong}}
			\scriptsize
			\begin{itemize}
				\item \wronganswer{High ranking on global GDP without AML laws} — This \textbf{sounds} like it would matter, but GDP alone \textbf{doesn't create secrecy}.
				\item \textbf{Minimal BO reporting:} If you don't know who owns companies, that's secrecy.
				\item \textbf{Cross-border services to non-residents:} Providing banking services to people who don't live there enables secrecy.
				\item \textbf{High secrecy-to-scale ratio:} How much secrecy relative to financial activity — if a small jurisdiction has a lot of secrecy, that's a red flag.
				\item \textbf{Key principle:} FSI measures \textbf{actual secrecy}, not just size or wealth.
			\end{itemize}
		\end{block}
		\vspace{0.02cm}
		\recallbox{\scriptsize \textbf{Key Rule:} FSI factors = \textbf{BO reporting + cross-border services + secrecy-to-scale ratio}.}
	\end{frame}
	
	% ================================================================
	% SLIDE 8: EWRA — THE MLRO ROLE TRAP
	% ================================================================
	\begin{frame}
		\frametitle{\fontsize{11}{13}\selectfont \textbf{EWRA — The MLRO Role Trap}}
		\begin{block}{\fontsize{7}{9}\selectfont \textbf{The Hard Question — CAMS Exam Favorite}}
			\scriptsize
			Which function uses EWRA results to shape risk posture?
		\end{block}
		\vspace{0.02cm}
		\answerbox{\large \textbf{The Money Laundering Reporting Officer (MLRO).}}
		\vspace{0.02cm}
		\begin{block}{\fontsize{7}{9}\selectfont \textbf{Natural Explanation — Why You Got This Wrong}}
			\scriptsize
			\begin{itemize}
				\item Many candidates choose \wronganswer{Internal Audit Lead} because audit \textbf{sounds} like it would shape risk.
				\item But \textbf{Internal Audit} is the \textbf{third line of defense} — they assess controls, they don't shape risk posture.
				\item The \textbf{MLRO} is responsible for \textbf{managing AML risk} — they use EWRA to adjust controls, enhance monitoring, and allocate resources.
				\item \textbf{Key principle:} MLRO = \textbf{risk manager}. Internal Audit = \textbf{risk assessor}. Different roles.
			\end{itemize}
		\end{block}
		\vspace{0.02cm}
		\recallbox{\scriptsize \textbf{Key Rule:} MLRO = \textbf{shapes risk posture}. Internal Audit = \textbf{assesses} controls.}
	\end{frame}
	
	% ================================================================
	% SLIDE 9: TAX EVASION vs AVOIDANCE — THE DISTINCTION TRAP
	% ================================================================
	\begin{frame}
		\frametitle{\fontsize{11}{13}\selectfont \textbf{Tax Evasion vs Tax Avoidance — The Distinction Trap}}
		\begin{block}{\fontsize{7}{9}\selectfont \textbf{The Hard Question — CAMS Exam Favorite}}
			\scriptsize
			A company failed to report offshore revenue and laundered funds through an unrelated third party. What conclusion is best supported?
		\end{block}
		\vspace{0.02cm}
		\answerbox{\large \textbf{The company committed tax evasion by concealing income.}}
		\vspace{0.02cm}
		\begin{block}{\fontsize{7}{9}\selectfont \textbf{Natural Explanation — Why You Got This Wrong}}
			\scriptsize
			\begin{itemize}
				\item Many candidates choose \wronganswer{legally avoided taxes through offshore structuring} because offshore structuring \textbf{sounds} like tax planning.
				\item But there's a \textbf{critical distinction}:
				\begin{itemize}
					\item \textbf{Tax avoidance:} Legal — you structure to minimize taxes \textbf{within the law}. Full disclosure required.
					\item \textbf{Tax evasion:} Illegal — you \textbf{conceal} income or assets. No disclosure.
				\end{itemize}
				\item \textbf{Failure to report + laundering through a third party} = concealment = \textbf{evasion}.
				\item \textbf{Key principle:} Tax avoidance requires \textbf{full disclosure}. Concealment = evasion = predicate offense for ML.
			\end{itemize}
		\end{block}
		\vspace{0.02cm}
		\recallbox{\scriptsize \textbf{Key Rule:} Avoidance = legal with disclosure. Evasion = illegal with concealment.}
	\end{frame}
	
	% ================================================================
	% SLIDE 10: BOARD OVERSIGHT — THE OPERATIONAL TRAP
	% ================================================================
	\begin{frame}
		\frametitle{\fontsize{11}{13}\selectfont \textbf{Board Oversight — The Operational Trap}}
		\begin{block}{\fontsize{7}{9}\selectfont \textbf{The Hard Question — CAMS Exam Favorite}}
			\scriptsize
			Which actions demonstrate effective board oversight? (Select Three.)
		\end{block}
		\vspace{0.02cm}
		\answerbox{\large \textbf{1. Receiving regular reports on risks. 2. Reviewing management's response to audit findings. 3. Challenging management on weak performance.}}
		\vspace{0.02cm}
		\begin{block}{\fontsize{7}{9}\selectfont \textbf{Natural Explanation — Why You Got This Wrong}}
			\scriptsize
			\begin{itemize}
				\item Many candidates choose \wronganswer{Participating in SAR investigations} because it sounds \textbf{engaged} and \textbf{involved}.
				\item But \textbf{boards provide OVERSIGHT}, not operational involvement.
				\item \textbf{Participating in SAR investigations} is \textbf{operational} — that's for compliance staff, not the board.
				\item \textbf{Board duties:}
				\begin{itemize}
					\item Receive reports (stay informed)
					\item Review audit responses (hold management accountable)
					\item Challenge management (ask tough questions)
				\end{itemize}
				\item \textbf{Key principle:} Board = \textbf{oversight and challenge}. NOT \textbf{doing the work}.
			\end{itemize}
		\end{block}
		\vspace{0.02cm}
		\recallbox{\scriptsize \textbf{Key Rule:} Board = \textbf{oversight}. NOT operational involvement.}
	\end{frame}
	
	% ================================================================
	% SLIDE 11: TERRORIST FINANCING — THE MOVEMENT METHODS TRAP
	% ================================================================
	\begin{frame}
		\frametitle{\fontsize{11}{13}\selectfont \textbf{Terrorist Financing — The Movement Methods Trap}}
		\begin{block}{\fontsize{7}{9}\selectfont \textbf{The Hard Question — CAMS Exam Favorite}}
			\scriptsize
			Which techniques are commonly used by terrorists to move or store money across borders avoiding detection? (Select Three.)
		\end{block}
		\vspace{0.02cm}
		\answerbox{\large \textbf{1. Fake charities. 2. Hawala (informal value transfer). 3. Physical cash transport across porous borders.}}
		\vspace{0.02cm}
		\begin{block}{\fontsize{7}{9}\selectfont \textbf{Natural Explanation — Why You Got This Wrong}}
			\scriptsize
			\begin{itemize}
				\item \wronganswer{Large wire transfers through regulated banks} — This is a \textbf{trap} because it \textbf{sounds} like a method, but it's \textbf{detectable}.
				\item Terrorists want to \textbf{avoid detection}, so they use \textbf{unregulated} methods.
				\item \textbf{Fake charities:} Disguise TF as humanitarian aid.
				\item \textbf{Hawala:} Informal value transfer outside regulated banking.
				\item \textbf{Physical cash transport:} Crosses borders avoiding detection.
				\item \textbf{Key principle:} Terrorists use \textbf{unregulated} and \textbf{hard-to-detect} methods.
			\end{itemize}
		\end{block}
		\vspace{0.02cm}
		\recallbox{\scriptsize \textbf{Key Rule:} TF = \textbf{fake charities + Hawala + physical cash}. NOT regulated wire transfers.}
	\end{frame}
	
	% ================================================================
	% SLIDE 12: HUMAN TRAFFICKING — THE LAUNDERING METHODS TRAP
	% ================================================================
	\begin{frame}
		\frametitle{\fontsize{11}{13}\selectfont \textbf{Human Trafficking — The Laundering Methods Trap}}
		\begin{block}{\fontsize{7}{9}\selectfont \textbf{The Hard Question — CAMS Exam Favorite}}
			\scriptsize
			Which characteristics enable human trafficking operations to effectively launder illicit proceeds? (Select Two.)
		\end{block}
		\vspace{0.02cm}
		\answerbox{\large \textbf{1. Use of mobile payment and remittance systems. 2. Weak cross-border enforcement cooperation.}}
		\vspace{0.02cm}
		\begin{block}{\fontsize{7}{9}\selectfont \textbf{Natural Explanation — Why You Got This Wrong}}
			\scriptsize
			\begin{itemize}
				\item \wronganswer{High correlation between trafficking and counterfeit document operations} — This is \textbf{related} to trafficking, but \textbf{counterfeit documents} are a \textbf{separate crime}, not a laundering method.
				\item \textbf{Mobile payments/remittances:} Enable rapid movement of small amounts (structuring).
				\item \textbf{Weak cross-border enforcement:} Allows traffickers to operate across jurisdictions.
				\item \textbf{Key principle:} Traffickers launder proceeds through \textbf{rapid, small, cross-border} transactions.
			\end{itemize}
		\end{block}
		\vspace{0.02cm}
		\recallbox{\scriptsize \textbf{Key Rule:} Trafficking laundering = \textbf{mobile payments} + \textbf{weak cross-border enforcement}.}
	\end{frame}
	
	% ================================================================
	% SLIDE 13: TCSP — THE OUTSOURCING TRAP
	% ================================================================
	\begin{frame}
		\frametitle{\fontsize{11}{13}\selectfont \textbf{TCSP — The Outsourcing Trap}}
		\begin{block}{\fontsize{7}{9}\selectfont \textbf{The Hard Question — CAMS Exam Favorite}}
			\scriptsize
			A TCSP is reviewing a client who wants confidentiality and is linked to high-risk industries in the NRA. What should the TCSP do?
		\end{block}
		\vspace{0.02cm}
		\answerbox{\large \textbf{Decline immediately unless full BO info and risk justifications are disclosed and verified.}}
		\vspace{0.02cm}
		\begin{block}{\fontsize{7}{9}\selectfont \textbf{Natural Explanation — Why You Got This Wrong}}
			\scriptsize
			\begin{itemize}
				\item \wronganswer{Refer to offshore local counsel and detach from compliance responsibility} — This is a \textbf{trap}. You \textbf{cannot} outsource AML responsibility.
				\item \textbf{TCSP retains compliance responsibility} — always.
				\item \textbf{Confidentiality requests + high-risk industries = red flags.}
				\item \textbf{FATF Rec 24/25:} BO identification is \textbf{mandatory}.
				\item \textbf{Key principle:} You can use local counsel, but you \textbf{cannot transfer liability}. You remain responsible.
			\end{itemize}
		\end{block}
		\vspace{0.02cm}
		\recallbox{\scriptsize \textbf{Key Rule:} TCSP cannot outsource AML responsibility. Decline unless BO disclosed.}
	\end{frame}
	
	% ================================================================
	% SLIDE 14: DNFBP RISK MANAGEMENT — THE EDD TRAP
	% ================================================================
	\begin{frame}
		\frametitle{\fontsize{11}{13}\selectfont \textbf{DNFBP Risk Management — The EDD Trap}}
		\begin{block}{\fontsize{7}{9}\selectfont \textbf{The Hard Question — CAMS Exam Favorite}}
			\scriptsize
			Which strategies help manage DNFBP risks? (Select Two.)
		\end{block}
		\vspace{0.02cm}
		\answerbox{\large \textbf{1. Creating a reputational risk/client selection committee. 2. Using standardized sector-specific questionnaires.}}
		\vspace{0.02cm}
		\begin{block}{\fontsize{7}{9}\selectfont \textbf{Natural Explanation — Why You Got This Wrong}}
			\scriptsize
			\begin{itemize}
				\item \wronganswer{Defaulting to EDD for all DNFBPs} — This \textbf{sounds} prudent, but it's \textbf{not risk-based}. You would over-control low-risk DNFBPs.
				\item \textbf{Risk-based approach:} Apply proportionate measures based on actual risk.
				\item \textbf{Client selection committee:} Governance for high-risk decisions.
				\item \textbf{Standardized questionnaires:} Consistent risk assessment across DNFBPs.
				\item \textbf{Key principle:} Risk management = \textbf{proportionate measures}, not blanket EDD.
			\end{itemize}
		\end{block}
		\vspace{0.02cm}
		\recallbox{\scriptsize \textbf{Key Rule:} DNFBP risk = \textbf{committee + standardized questionnaires}. NOT blanket EDD.}
	\end{frame}
	
	% ================================================================
	% SLIDE 15: LAYERING — THE CLASSIFICATION TRAP
	% ================================================================
	\begin{frame}
		\frametitle{\fontsize{11}{13}\selectfont \textbf{Layering — The Classification Trap}}
		\begin{block}{\fontsize{7}{9}\selectfont \textbf{The Hard Question — CAMS Exam Favorite}}
			\scriptsize
			A private loan extended to a shell company with no operations, loan amount matches criminal proceeds. What is the initial classification?
		\end{block}
		\vspace{0.02cm}
		\answerbox{\large \textbf{A layering tactic designed to conceal the origin of illicit funds.}}
		\vspace{0.02cm}
		\begin{block}{\fontsize{7}{9}\selectfont \textbf{Natural Explanation — Why You Got This Wrong}}
			\scriptsize
			\begin{itemize}
				\item \wronganswer{A placement strategy to integrate clean capital} — This is a \textbf{trap}. Placement is the \textbf{first stage} — entering illicit funds into the system.
				\item \textbf{Layering} is the \textbf{second stage} — concealing the origin.
				\item A \textbf{private loan to a shell company} creates a \textbf{false paper trail} that appears legitimate.
				\item The loan \textbf{distances} the funds from their criminal source.
				\item \textbf{Key principle:} Layering = \textbf{concealing origin}. Placement = \textbf{entering system}. Integration = \textbf{using funds legitimately}.
			\end{itemize}
		\end{block}
		\vspace{0.02cm}
		\recallbox{\scriptsize \textbf{Key Rule:} Layering = \textbf{concealing origin}. Placement = \textbf{entering system}.}
	\end{frame}
	
	% ================================================================
	% SLIDE 16: PROACTIVE RESPONSE — THE UPDATING TRAP
	% ================================================================
	\begin{frame}
		\frametitle{\fontsize{11}{13}\selectfont \textbf{Proactive Response — The Updating Trap}}
		\begin{block}{\fontsize{7}{9}\selectfont \textbf{The Hard Question — CAMS Exam Favorite}}
			\scriptsize
			What is a proactive institutional response to newly detected AML risks?
		\end{block}
		\vspace{0.02cm}
		\answerbox{\large \textbf{Immediately updating the enterprise-wide risk assessment.}}
		\vspace{0.02cm}
		\begin{block}{\fontsize{7}{9}\selectfont \textbf{Natural Explanation — Why You Got This Wrong}}
			\scriptsize
			\begin{itemize}
				\item \wronganswer{Reducing reliance on transaction monitoring systems} — This is \textbf{reactive} and \textbf{counterproductive}.
				\item \textbf{Proactive} = take action \textbf{before} the risk materializes.
				\item \textbf{Updating EWRA} allows the institution to:
				\begin{itemize}
					\item Adjust controls
					\item Allocate resources to higher-risk areas
					\item Enhance monitoring before risks materialize
				\end{itemize}
				\item \textbf{Key principle:} Proactive = \textbf{update risk assessment}. Reactive = \textbf{respond after the fact}.
			\end{itemize}
		\end{block}
		\vspace{0.02cm}
		\recallbox{\scriptsize \textbf{Key Rule:} Proactive = \textbf{update EWRA immediately}.}
	\end{frame}
	
	% ================================================================
	% SLIDE 17: IOSCO — THE OBJECTIVES TRAP
	% ================================================================
	\begin{frame}
		\frametitle{\fontsize{11}{13}\selectfont \textbf{IOSCO — The Objectives Trap}}
		\begin{block}{\fontsize{7}{9}\selectfont \textbf{The Hard Question — CAMS Exam Favorite}}
			\scriptsize
			Which of IOSCO's core objectives directly support anti-financial crime efforts? (Select Two.)
		\end{block}
		\vspace{0.02cm}
		\answerbox{\large \textbf{1. Enhancing investor protection through cooperation and enforcement. 2. Promoting financial stability by managing systemic risks and facilitating international information exchange.}}
		\vspace{0.02cm}
		\begin{block}{\fontsize{7}{9}\selectfont \textbf{Natural Explanation — Why You Got This Wrong}}
			\scriptsize
			\begin{itemize}
				\item \wronganswer{Requiring uniform compliance departments} — This \textbf{sounds} like it would help, but it's \textbf{not} an IOSCO objective.
				\item \textbf{IOSCO objectives:}
				\begin{itemize}
					\item \textbf{Investor protection:} Combating fraud and market abuse.
					\item \textbf{Financial stability:} Managing systemic risks and cross-border info sharing.
				\end{itemize}
				\item \textbf{Key principle:} IOSCO's objectives are \textbf{high-level} — investor protection and financial stability. Uniform departments is \textbf{operational}, not an objective.
			\end{itemize}
		\end{block}
		\vspace{0.02cm}
		\recallbox{\scriptsize \textbf{Key Rule:} IOSCO = \textbf{investor protection + financial stability}.}
	\end{frame}
	
	% ================================================================
	% SLIDE 18: IMF — ECONOMIC SUCCESS vs AML EFFECTIVENESS
	% ================================================================
	\begin{frame}
		\frametitle{\fontsize{11}{13}\selectfont \textbf{IMF — The Economic Success Trap}}
		\begin{block}{\fontsize{7}{9}\selectfont \textbf{The Hard Question — CAMS Exam Favorite}}
			\scriptsize
			A nation achieved strong GDP growth after IMF reforms, but the IMF still expresses concerns. Which explanation best reflects the IMF's approach?
		\end{block}
		\vspace{0.02cm}
		\answerbox{\large \textbf{Improvements in economic measures can mask financial crime vulnerabilities.}}
		\vspace{0.02cm}
		\begin{block}{\fontsize{7}{9}\selectfont \textbf{Natural Explanation — Why You Got This Wrong}}
			\scriptsize
			\begin{itemize}
				\item Many candidates think: \wronganswer{The IMF is being unfair} or \wronganswer{The country needs more reforms}.
				\item \textbf{Key insight:} Strong GDP growth \textbf{does not} mean strong AML controls.
				\item A country can have \textbf{economic success} but \textbf{weak AML frameworks}.
				\item \textbf{IMF assessment looks BEYOND economic indicators} to evaluate financial crime vulnerabilities.
				\item \textbf{Key principle:} Economic success ≠ AML effectiveness. IMF evaluates \textbf{financial crime vulnerabilities} separately.
			\end{itemize}
		\end{block}
		\vspace{0.02cm}
		\recallbox{\scriptsize \textbf{Key Rule:} Economic success ≠ AML effectiveness. IMF looks beyond GDP.}
	\end{frame}
	
	% ================================================================
	% SLIDE 19: OPEN SOURCE RESEARCH — THE ENTITY-SPECIFIC TRAP
	% ================================================================
	\begin{frame}
		\frametitle{\fontsize{11}{13}\selectfont \textbf{Open Source Research — The Entity-Specific Trap}}
		\begin{block}{\fontsize{7}{9}\selectfont \textbf{The Hard Question — CAMS Exam Favorite}}
			\scriptsize
			An analyst notes wire transfers to an unfamiliar entity in a sanctioned jurisdiction. Which information source would most help understand the activity?
		\end{block}
		\vspace{0.02cm}
		\answerbox{\large \textbf{Open-source research into the entity.}}
		\vspace{0.02cm}
		\begin{block}{\fontsize{7}{9}\selectfont \textbf{Natural Explanation — Why You Got This Wrong}}
			\scriptsize
			\begin{itemize}
				\item \wronganswer{Regulatory advisories about common evasion techniques} — This is \textbf{general} information, not \textbf{specific} to the entity.
				\item \textbf{Open-source research} on the \textbf{specific entity} provides:
				\begin{itemize}
					\item Ownership information
					\item Business activities
					\item Connections to other entities
					\item Whether it's a shell company or front
				\end{itemize}
				\item \textbf{Key principle:} When investigating an \textbf{unknown entity}, go to \textbf{entity-specific} sources, not general guidance.
			\end{itemize}
		\end{block}
		\vspace{0.02cm}
		\recallbox{\scriptsize \textbf{Key Rule:} Unknown entity = \textbf{open-source research} on that specific entity.}
	\end{frame}
	
	% ================================================================
	% SLIDE 20: BATCH SCREENING — THE VOLUME TRAP
	% ================================================================
	\begin{frame}
		\frametitle{\fontsize{11}{13}\selectfont \textbf{Batch Screening — The Volume Trap}}
		\begin{block}{\fontsize{7}{9}\selectfont \textbf{The Hard Question — CAMS Exam Favorite}}
			\scriptsize
			Which are true about batch screening? (Select Two.)
		\end{block}
		\vspace{0.02cm}
		\answerbox{\large \textbf{1. It helps reveal changing risks within the existing customer base. 2. It is designed to handle large volumes of data.}}
		\vspace{0.02cm}
		\begin{block}{\fontsize{7}{9}\selectfont \textbf{Natural Explanation — Why You Got This Wrong}}
			\scriptsize
			\begin{itemize}
				\item \wronganswer{It provides real-time transaction screening} — This is \textbf{real-time screening}, not batch screening.
				\item \textbf{Batch screening:} Processes data in bulk — typically the entire customer database — against sanctions lists, PEP lists.
				\item \textbf{Purpose:} Detects existing customers who have become risky (e.g., newly designated sanctions).
				\item \textbf{Volume-oriented:} Designed to handle millions of records.
				\item \textbf{Key principle:} Batch screening = \textbf{periodic re-screening} of the \textbf{entire customer base}.
			\end{itemize}
		\end{block}
		\vspace{0.02cm}
		\recallbox{\scriptsize \textbf{Key Rule:} Batch screening = \textbf{large volumes + changing risks}. NOT real-time.}
	\end{frame}
	
	% ================================================================
	% SLIDE 21: DATA EXTRACTION — THE NORMALIZATION TRAP
	% ================================================================
	\begin{frame}
		\frametitle{\fontsize{11}{13}\selectfont \textbf{Data Extraction — The Normalization Trap}}
		\begin{block}{\fontsize{7}{9}\selectfont \textbf{The Hard Question — CAMS Exam Favorite}}
			\scriptsize
			Which best describes a challenge during AFC data extraction?
		\end{block}
		\vspace{0.02cm}
		\answerbox{\large \textbf{Data may come in different formats, requiring normalization during integration.}}
		\vspace{0.02cm}
		\begin{block}{\fontsize{7}{9}\selectfont \textbf{Natural Explanation — Why You Got This Wrong}}
			\scriptsize
			\begin{itemize}
				\item Many candidates think data extraction is just \textbf{pulling data} from systems. But the \textbf{real challenge} is what happens \textbf{after} you pull it.
				\item Different systems use \textbf{different formats}: one system uses "US," another uses "USA," another uses "840" or "United States."
				\item \textbf{Normalization} is the process of making all this data \textbf{consistent} so it can be used together.
				\item Without normalization, your transaction monitoring system can't properly analyze transactions or identify patterns.
				\item \textbf{Key principle:} Data extraction is \textbf{not} the challenge — \textbf{data normalization} is.
			\end{itemize}
		\end{block}
		\vspace{0.02cm}
		\recallbox{\scriptsize \textbf{Key Rule:} Data extraction challenge = \textbf{different formats requiring normalization}.}
	\end{frame}
	
	% ================================================================
	% SLIDE 22: DATA MINING — THE EXTERNAL SOURCES TRAP
	% ================================================================
	\begin{frame}
		\frametitle{\fontsize{11}{13}\selectfont \textbf{Data Mining — The External Sources Trap}}
		\begin{block}{\fontsize{7}{9}\selectfont \textbf{The Hard Question — CAMS Exam Favorite}}
			\scriptsize
			Which are requirements for effective data mining and matching in AFC controls? (Select Two.)
		\end{block}
		\vspace{0.02cm}
		\answerbox{\large \textbf{1. Comparing extracted data against structured and unstructured risk sources. 2. Applying intelligent filters to reduce false positives during entity resolution.}}
		\vspace{0.02cm}
		\begin{block}{\fontsize{7}{9}\selectfont \textbf{Natural Explanation — Why You Got This Wrong}}
			\scriptsize
			\begin{itemize}
				\item \wronganswer{Focusing on internal customer files for mining purposes} — This is a \textbf{trap}. Effective data mining requires \textbf{both internal and external} sources.
				\item \textbf{Structured data:} Databases, spreadsheets, formatted fields.
				\item \textbf{Unstructured data:} Emails, documents, free-text fields, news articles.
				\item \textbf{Intelligent filters:} Reduce false positives during entity resolution.
				\item \textbf{Key principle:} Data mining is \textbf{broad}, not narrow. You need \textbf{external} data to catch \textbf{external} risks.
			\end{itemize}
		\end{block}
		\vspace{0.02cm}
		\recallbox{\scriptsize \textbf{Key Rule:} Data mining = \textbf{internal + external} data sources + \textbf{intelligent filters}.}
	\end{frame}
	
	% ================================================================
	% SLIDE 23: DATA INTEGRATION — THE ISOLATION TRAP
	% ================================================================
	\begin{frame}
		\frametitle{\fontsize{11}{13}\selectfont \textbf{Data Integration — The Isolation Trap}}
		\begin{block}{\fontsize{7}{9}\selectfont \textbf{The Hard Question — CAMS Exam Favorite}}
			\scriptsize
			Which help ensure effective integration of multiple data sources into compliance platforms? (Select Three.)
		\end{block}
		\vspace{0.02cm}
		\answerbox{\large \textbf{1. Maintaining a common data dictionary. 2. Mapping all incoming data to standardized field names. 3. Aligning control data requirements to system inputs.}}
		\vspace{0.02cm}
		\begin{block}{\fontsize{7}{9}\selectfont \textbf{Natural Explanation — Why You Got This Wrong}}
			\scriptsize
			\begin{itemize}
				\item \wronganswer{Separating internal and external data into isolated systems} — This is a \textbf{trap}. Separation \textbf{undermines} integration, it doesn't help it.
				\item \textbf{Common data dictionary:} Ensures all systems use the same definitions.
				\item \textbf{Field mapping:} Ensures data from source systems is translated correctly.
				\item \textbf{Aligning control requirements:} Ensures the monitoring system receives all necessary data.
				\item \textbf{Key principle:} Integration is about \textbf{connection}, not isolation.
			\end{itemize}
		\end{block}
		\vspace{0.02cm}
		\recallbox{\scriptsize \textbf{Key Rule:} Integration = \textbf{common dictionary + field mapping + aligned requirements}.}
	\end{frame}
	
	% ================================================================
	% SLIDE 24: INTERNAL vs EXTERNAL SOURCES — THE SCOPE TRAP
	% ================================================================
	\begin{frame}
		\frametitle{\fontsize{11}{13}\selectfont \textbf{Internal vs External Sources — The Scope Trap}}
		\begin{block}{\fontsize{7}{9}\selectfont \textbf{The Hard Question — CAMS Exam Favorite}}
			\scriptsize
			Which characteristics describe internal versus external sources used during AFC data extraction? (Select Two.)
		\end{block}
		\vspace{0.02cm}
		\answerbox{\large \textbf{1. Internal sources reflect observed customer behavior or system-held values. 2. External sources include sanction or adverse media lists from agencies.}}
		\vspace{0.02cm}
		\begin{block}{\fontsize{7}{9}\selectfont \textbf{Natural Explanation — Why You Got This Wrong}}
			\scriptsize
			\begin{itemize}
				\item \wronganswer{Internal data includes account balance snapshots, not user activity} — This is \textbf{too narrow}. Internal data includes \textbf{much more} than just account balances.
				\item \textbf{Internal sources:} Transaction histories, account balances, customer profiles, onboarding records.
				\item \textbf{External sources:} Sanctions lists (OFAC, UN), PEP lists, adverse media databases.
				\item \textbf{Key principle:} Internal data provides \textbf{behavioral context}. External data provides \textbf{risk context}. Both are essential.
			\end{itemize}
		\end{block}
		\vspace{0.02cm}
		\recallbox{\scriptsize \textbf{Key Rule:} Internal = \textbf{behavioral context}. External = \textbf{risk context}. Both needed.}
	\end{frame}
	
	% ================================================================
	% SLIDE 25: BSA POLICY — THE TAILORING TRAP
	% ================================================================
	\begin{frame}
		\frametitle{\fontsize{11}{13}\selectfont \textbf{BSA Policy — The Tailoring Trap}}
		\begin{block}{\fontsize{7}{9}\selectfont \textbf{The Hard Question — CAMS Exam Favorite}}
			\scriptsize
			Which are key features of well-structured AML/CFT policies aligned with the BSA? (Select Two.)
		\end{block}
		\vspace{0.02cm}
		\answerbox{\large \textbf{1. Tailoring procedures to specific business units or jurisdictions. 2. Assigning roles and responsibilities to specific individuals or teams.}}
		\vspace{0.02cm}
		\begin{block}{\fontsize{7}{9}\selectfont \textbf{Natural Explanation — Why You Got This Wrong}}
			\scriptsize
			\begin{itemize}
				\item \wronganswer{Requiring transaction reporting above a fixed institutional threshold} — This is \textbf{not risk-based}. Fixed thresholds ignore risk differences.
				\item \textbf{Tailoring:} Different business units have different risk profiles — procedures should reflect that.
				\item \textbf{Assigning roles:} Accountability is critical — someone must be responsible.
				\item \textbf{Key principle:} BSA requires a \textbf{risk-based} approach, not a one-size-fits-all approach.
			\end{itemize}
		\end{block}
		\vspace{0.02cm}
		\recallbox{\scriptsize \textbf{Key Rule:} BSA policies = \textbf{tailoring + assigning roles}. NOT fixed thresholds.}
	\end{frame}
	
	% ================================================================
	% SLIDE 26: CHINA AML LAW 2025 — THE EXPANSION TRAP
	% ================================================================
	\begin{frame}
		\frametitle{\fontsize{11}{13}\selectfont \textbf{China AML Law 2025 — The Expansion Trap}}
		\begin{block}{\fontsize{7}{9}\selectfont \textbf{The Hard Question — CAMS Exam Favorite}}
			\scriptsize
			What is an accurate statement regarding China's AML Law which took effect in 2025?
		\end{block}
		\vspace{0.02cm}
		\answerbox{\large \textbf{It added coverage for additional sectors including law firms, real estate agencies, and dealers in precious gems.}}
		\vspace{0.02cm}
		\begin{block}{\fontsize{7}{9}\selectfont \textbf{Natural Explanation — Why You Got This Wrong}}
			\scriptsize
			\begin{itemize}
				\item \wronganswer{It reduced coverage to only banks} — \textbf{NO}. It \textbf{expanded} coverage.
				\item \wronganswer{It eliminated all reporting requirements} — \textbf{NO}. It strengthened them.
				\item \textbf{China's 2025 AML Law} added \textbf{DNFBPs} (designated non-financial businesses and professions) to its scope.
				\item This aligns China with \textbf{FATF standards} requiring DNFBPs to implement AML/CFT obligations.
				\item \textbf{Key principle:} The trend is \textbf{expansion}, not reduction.
			\end{itemize}
		\end{block}
		\vspace{0.02cm}
		\recallbox{\scriptsize \textbf{Key Rule:} China AML 2025 = \textbf{expanded coverage} to DNFBPs.}
	\end{frame}
	
	% ================================================================
	% SLIDE 27: EU AMLA — THE CROSS-BORDER TRAP
	% ================================================================
	\begin{frame}
		\frametitle{\fontsize{11}{13}\selectfont \textbf{EU AMLA — The Cross-Border Trap}}
		\begin{block}{\fontsize{7}{9}\selectfont \textbf{The Scenario — CAMS Exam Favorite}}
			\scriptsize
			An EU-based fintech operates in \textbf{5 member states} and has been flagged for non-compliance by \textbf{more than one regulator}.
		\end{block}
		\vspace{0.02cm}
		\begin{block}{\fontsize{7}{9}\selectfont \textbf{The Hard Question}}
			\scriptsize
			Which supervisory authority is most likely to assume oversight?
		\end{block}
		\vspace{0.02cm}
		\answerbox{\large \textbf{The AMLA, due to the firm's cross-border nature and risk profile.}}
		\vspace{0.02cm}
		\begin{block}{\fontsize{7}{9}\selectfont \textbf{Natural Explanation — Why You Got This Wrong}}
			\scriptsize
			\begin{itemize}
				\item \wronganswer{The national regulator of headquarters} — This is \textbf{traditional} supervision, but AMLA changes this.
				\item \textbf{AMLA (Anti-Money Laundering Authority)} = new EU body (established 2024).
				\item \textbf{Direct supervision:} AMLA directly supervises selected \textbf{cross-border financial institutions}.
				\item Flagged by \textbf{multiple regulators} = cross-border risk = AMLA jurisdiction.
				\item \textbf{Key principle:} AMLA = \textbf{centralized supervision} for cross-border, high-risk EU institutions.
			\end{itemize}
		\end{block}
		\vspace{0.02cm}
		\recallbox{\scriptsize \textbf{Key Rule:} AMLA = \textbf{centralized supervision} for cross-border, high-risk EU institutions.}
	\end{frame}
	
	% ================================================================
	% SLIDE 28: UK RISK-BASED APPROACH — THE PURPOSE TRAP
	% ================================================================
	\begin{frame}
		\frametitle{\fontsize{11}{13}\selectfont \textbf{UK Risk-Based Approach — The Purpose Trap}}
		\begin{block}{\fontsize{7}{9}\selectfont \textbf{The Hard Question — CAMS Exam Favorite}}
			\scriptsize
			Under the UK's AML framework, what is the primary purpose of the risk-based approach?
		\end{block}
		\vspace{0.02cm}
		\answerbox{\large \textbf{To allow firms to customize controls suited to their risk exposure.}}
		\vspace{0.02cm}
		\begin{block}{\fontsize{7}{9}\selectfont \textbf{Natural Explanation — Why You Got This Wrong}}
			\scriptsize
			\begin{itemize}
				\item \wronganswer{To eliminate all AML controls} — \textbf{NO.} Controls are still required.
				\item \wronganswer{To require the same controls for all firms} — That's \textbf{prescriptive}, not risk-based.
				\item \textbf{Customize controls to risk} = allocate resources to highest risk areas.
				\item \textbf{NOT} zero due diligence — some CDD always required.
				\item \textbf{NOT} de-risking — cannot simply decline entire categories.
				\item \textbf{Key principle:} Risk-based approach = \textbf{customize controls} to risk exposure.
			\end{itemize}
		\end{block}
		\vspace{0.02cm}
		\recallbox{\scriptsize \textbf{Key Rule:} Risk-based approach = \textbf{customize controls to risk}.}
	\end{frame}
	
	% ================================================================
	% SLIDE 29: CONCENTRATION ACCOUNT — THE TRACEABILITY TRAP
	% ================================================================
	\begin{frame}
		\frametitle{\fontsize{11}{13}\selectfont \textbf{Concentration Account — The Traceability Trap}}
		\begin{block}{\fontsize{7}{9}\selectfont \textbf{The Hard Question — CAMS Exam Favorite}}
			\scriptsize
			A treasury unit uses a concentration account lacking sufficient oversight. What is the most likely regulatory concern?
		\end{block}
		\vspace{0.02cm}
		\answerbox{\large \textbf{The institution might inadvertently facilitate money laundering due to poor traceability.}}
		\vspace{0.02cm}
		\begin{block}{\fontsize{7}{9}\selectfont \textbf{Natural Explanation — Why You Got This Wrong}}
			\scriptsize
			\begin{itemize}
				\item \wronganswer{The account has high transaction fees} — Not a regulatory concern.
				\item \wronganswer{The account violates tax laws} — Not the primary concern.
				\item \textbf{Concentration accounts (nostro/vostro)} aggregate funds from multiple customers.
				\item Without oversight, it's difficult to identify the \textbf{ultimate originator or beneficiary}.
				\item This \textbf{inadvertently facilitates money laundering} by obscuring the source of funds.
				\item \textbf{Key principle:} Concentration accounts require \textbf{strong oversight} for traceability.
			\end{itemize}
		\end{block}
		\vspace{0.02cm}
		\recallbox{\scriptsize \textbf{Key Rule:} Concentration account risk = \textbf{poor traceability} = ML facilitation.}
	\end{frame}
	
	% ================================================================
	% SLIDE 30: TRADE TOPOLOGY — THE PENNY STOCK TRAP
	% ================================================================
	\begin{frame}
		\frametitle{\fontsize{11}{13}\selectfont \textbf{Trade Topology — The Penny Stock Trap}}
		\begin{block}{\fontsize{7}{9}\selectfont \textbf{The Hard Question — CAMS Exam Favorite}}
			\scriptsize
			Two trades: high-volume equities with normal volatility vs low-volume penny stocks with rapid pricing shifts. Both passed basic monitoring. Which should be prioritized?
		\end{block}
		\vspace{0.02cm}
		\answerbox{\large \textbf{The penny stock trade, due to higher manipulation risk.}}
		\vspace{0.02cm}
		\begin{block}{\fontsize{7}{9}\selectfont \textbf{Natural Explanation — Why You Got This Wrong}}
			\scriptsize
			\begin{itemize}
				\item \wronganswer{Both equally, as volume and volatility are neutral indicators} — This ignores \textbf{risk-based} prioritization.
				\item \textbf{Low-volume penny stocks} are more susceptible to \textbf{manipulation} (pump-and-dump, wash trading).
				\item \textbf{High-volume equities} with normal volatility are less suspicious.
				\item Passing basic monitoring \textbf{does not} eliminate the need for risk-based prioritization.
				\item \textbf{Key principle:} Risk-based approach = prioritize \textbf{higher-risk} activities.
			\end{itemize}
		\end{block}
		\vspace{0.02cm}
		\recallbox{\scriptsize \textbf{Key Rule:} Penny stocks = \textbf{higher manipulation risk}. Prioritize them.}
	\end{frame}
	
	% ================================================================
	% SLIDE 31: RISK-BASED AUTHENTICATION — THE USER EXPERIENCE TRAP
	% ================================================================
	\begin{frame}
		\frametitle{\fontsize{11}{13}\selectfont \textbf{Risk-Based Authentication — The UX Trap}}
		\begin{block}{\fontsize{7}{9}\selectfont \textbf{The Hard Question — CAMS Exam Favorite}}
			\scriptsize
			How do risk-based authentication systems optimize user experience?
		\end{block}
		\vspace{0.02cm}
		\answerbox{\large \textbf{They prompt additional checks only when user behavior deviates from expectations.}}
		\vspace{0.02cm}
		\begin{block}{\fontsize{7}{9}\selectfont \textbf{Natural Explanation — Why You Got This Wrong}}
			\scriptsize
			\begin{itemize}
				\item \wronganswer{They eliminate all authentication requirements} — \textbf{NO.} Security still required.
				\item \wronganswer{They always require MFA for all users} — That's \textbf{not risk-based}, it's \textbf{uniform}.
				\item \textbf{Risk-based authentication (RBA)} uses \textbf{behavioral analytics} to establish a baseline of "normal" behavior.
				\item When behavior deviates, RBA prompts additional checks.
				\item Low-risk users have \textbf{frictionless} experiences; high-risk activities trigger additional authentication.
				\item \textbf{Key principle:} RBA balances \textbf{security and convenience}.
			\end{itemize}
		\end{block}
		\vspace{0.02cm}
		\recallbox{\scriptsize \textbf{Key Rule:} RBA = \textbf{prompt only when behavior deviates}.}
	\end{frame}
	
	% ================================================================
	% SLIDE 32: OFAC — THE SANCTIONS TRAP
	% ================================================================
	\begin{frame}
		\frametitle{\fontsize{11}{13}\selectfont \textbf{OFAC — The Sanctions Trap}}
		\begin{block}{\fontsize{7}{9}\selectfont \textbf{The Hard Question — CAMS Exam Favorite}}
			\scriptsize
			What is the primary focus of OFAC in U.S. financial regulation?
		\end{block}
		\vspace{0.02cm}
		\answerbox{\large \textbf{Issuing and enforcing economic and trade sanctions.}}
		\vspace{0.02cm}
		\begin{block}{\fontsize{7}{9}\selectfont \textbf{Natural Explanation — Why You Got This Wrong}}
			\scriptsize
			\begin{itemize}
				\item \wronganswer{Regulating anti-money laundering programs} — That's \textbf{FinCEN}, not OFAC.
				\item \wronganswer{Supervising banks} — That's \textbf{Federal Reserve/OCC}, not OFAC.
				\item \textbf{OFAC = sanctions enforcement.}
				\item \textbf{SDN List:} Specially Designated Nationals and Blocked Persons.
				\item \textbf{50\% Rule:} If sanctioned entity owns 50\%+ of another entity, that entity is also blocked.
				\item \textbf{Strict liability:} No intent required.
				\item \textbf{Key principle:} OFAC = sanctions. FinCEN = AML/BSA. Don't confuse them.
			\end{itemize}
		\end{block}
		\vspace{0.02cm}
		\recallbox{\scriptsize \textbf{Key Rule:} OFAC = \textbf{sanctions}. FinCEN = \textbf{AML/BSA}.}
	\end{frame}
	
	% ================================================================
	% SLIDE 33: E-COMMERCE MITIGATION — THE TIERED ONBOARDING TRAP
	% ================================================================
	\begin{frame}
		\frametitle{\fontsize{11}{13}\selectfont \textbf{E-Commerce Mitigation — The Tiered Onboarding Trap}}
		\begin{block}{\fontsize{7}{9}\selectfont \textbf{The Hard Question — CAMS Exam Favorite}}
			\scriptsize
			Which strategies help e-commerce platforms mitigate ML risks? (Select Two.)
		\end{block}
		\vspace{0.02cm}
		\answerbox{\large \textbf{1. Introducing tiered onboarding based on projected sales activity. 2. Conducting sanctions screening on all sellers during onboarding.}}
		\vspace{0.02cm}
		\begin{block}{\fontsize{7}{9}\selectfont \textbf{Natural Explanation — Why You Got This Wrong}}
			\scriptsize
			\begin{itemize}
				\item \wronganswer{Implementing seller-rankings based on transaction volume} — This could \textbf{incentivize ML activity} by encouraging high-volume sellers to launder more.
				\item \textbf{Tiered onboarding:} Scales CDD to risk level (higher volume = more scrutiny).
				\item \textbf{Sanctions screening:} Catches sanctioned entities before they join the platform.
				\item \textbf{Key principle:} Mitigation = \textbf{risk-based CDD + sanctions screening}.
			\end{itemize}
		\end{block}
		\vspace{0.02cm}
		\recallbox{\scriptsize \textbf{Key Rule:} E-commerce mitigation = \textbf{tiered onboarding + sanctions screening}.}
	\end{frame}
	
	% ================================================================
	% SLIDE 34: TOKENS — DECENTRALIZED vs CENTRALIZED TRAP
	% ================================================================
	\begin{frame}
		\frametitle{\fontsize{11}{13}\selectfont \textbf{Tokens — The Decentralized vs Centralized Trap}}
		\begin{block}{\fontsize{7}{9}\selectfont \textbf{The Hard Question — CAMS Exam Favorite}}
			\scriptsize
			Which are accurate statements regarding tokens in ML activities? (Select Two.)
		\end{block}
		\vspace{0.02cm}
		\answerbox{\large \textbf{Tokens may be exchanged rapidly on decentralized platforms, obscuring source and destination of funds.}}
		\vspace{0.02cm}
		\begin{block}{\fontsize{7}{9}\selectfont \textbf{Natural Explanation — Why You Got This Wrong}}
			\scriptsize
			\begin{itemize}
				\item \wronganswer{Tokens may be exchanged rapidly on centralized platforms, obscuring source and destination} — \textbf{NO.} Centralized platforms require KYC and provide tracing.
				\item \textbf{Decentralized platforms (DEXs)} allow rapid token exchanges \textbf{without KYC}.
				\item This \textbf{obscures} source and destination of funds.
				\item \textbf{Centralized platforms} require KYC and are traceable.
				\item \textbf{Key principle:} Decentralized = anonymity. Centralized = KYC and traceability.
			\end{itemize}
		\end{block}
		\vspace{0.02cm}
		\recallbox{\scriptsize \textbf{Key Rule:} Decentralized = anonymity. Centralized = KYC and traceability.}
	\end{frame}
	
	% ================================================================
	% SLIDE 35: HIGH-VALUE ASSETS — THE CRYPTO RED FLAG TRAP
	% ================================================================
	\begin{frame}
		\frametitle{\fontsize{11}{13}\selectfont \textbf{High-Value Assets — The Crypto Red Flag Trap}}
		\begin{block}{\fontsize{7}{9}\selectfont \textbf{The Hard Question — CAMS Exam Favorite}}
			\scriptsize
			Which red flag is most commonly associated with abuse of high-value assets for laundering?
		\end{block}
		\vspace{0.02cm}
		\answerbox{\large \textbf{Making payments for high-value items in cryptocurrency.}}
		\vspace{0.02cm}
		\begin{block}{\fontsize{7}{9}\selectfont \textbf{Natural Explanation — Why You Got This Wrong}}
			\scriptsize
			\begin{itemize}
				\item \wronganswer{Purchasing assets with complex financing plans} — This \textbf{can} be a red flag, but crypto is \textbf{more directly} associated.
				\item \textbf{Cryptocurrency offers pseudonymity}, making it attractive for illicit payments.
				\item \textbf{High-value items:} Real estate, art, luxury cars purchased with crypto = suspicion.
				\item \textbf{Key principle:} Crypto payments for high-value assets = \textbf{primary red flag}.
			\end{itemize}
		\end{block}
		\vspace{0.02cm}
		\recallbox{\scriptsize \textbf{Key Rule:} Crypto payments for high-value assets = \textbf{primary red flag}.}
	\end{frame}
	
	% ================================================================
	% SLIDE 36: PSP — THE CUSTOMIZATION TRAP
	% ================================================================
	\begin{frame}
		\frametitle{\fontsize{11}{13}\selectfont \textbf{PSP — The Customization Trap}}
		\begin{block}{\fontsize{7}{9}\selectfont \textbf{The Scenario — CAMS Exam Favorite}}
			\scriptsize
			A PSP is planning to offer new services targeting small merchants in multiple countries, some lacking mature AML frameworks.
		\end{block}
		\vspace{0.02cm}
		\begin{block}{\fontsize{7}{9}\selectfont \textbf{The Hard Question}}
			\scriptsize
			Which strategy is most appropriate?
		\end{block}
		\vspace{0.02cm}
		\answerbox{\large \textbf{Customize AML policy and onboarding procedures per jurisdiction and product.}}
		\vspace{0.02cm}
		\begin{block}{\fontsize{7}{9}\selectfont \textbf{Natural Explanation — Why You Got This Wrong}}
			\scriptsize
			\begin{itemize}
				\item \wronganswer{Apply home-country AML standards across all jurisdictions} — \textbf{NO.} One-size-fits-all fails.
				\item \textbf{Jurisdiction-specific:} High-risk = EDD, transaction limits, or restricted services.
				\item \textbf{Product-specific:} Low-value vs high-value = different risk profiles.
				\item \textbf{Key principle:} PSPs must \textbf{customize} AML per jurisdiction and product.
			\end{itemize}
		\end{block}
		\vspace{0.02cm}
		\recallbox{\scriptsize \textbf{Key Rule:} PSPs = \textbf{customize per jurisdiction and product}.}
	\end{frame}
	
	% ================================================================
	% SLIDE 37: CASINO CONTROLS — THE SELECTION TRAP
	% ================================================================
	\begin{frame}
		\frametitle{\fontsize{11}{13}\selectfont \textbf{Casino Controls — The Selection Trap}}
		\begin{block}{\fontsize{7}{9}\selectfont \textbf{The Hard Question — CAMS Exam Favorite}}
			\scriptsize
			Which control measures help mitigate AML risks in casinos? (Select Two.)
		\end{block}
		\vspace{0.02cm}
		\answerbox{\large \textbf{1. Enforcing chip redemption policies tied to minimum gaming thresholds. 2. Tracking player activity through registered loyalty programs.}}
		\vspace{0.02cm}
		\begin{block}{\fontsize{7}{9}\selectfont \textbf{Natural Explanation — Why You Got This Wrong}}
			\scriptsize
			\begin{itemize}
				\item \wronganswer{Monitoring unstructured table gameplay for win/loss patterns} — This is a \textbf{monitoring technique}, not a formal control measure.
				\item \textbf{Chip redemption policies:} Prevents "chip washing" — exchanging chips for cash without meaningful gaming.
				\item \textbf{Loyalty programs:} Provide rich data on player behavior (wins, losses, session durations).
				\item \textbf{Key principle:} Controls = \textbf{formal policies + tracking mechanisms}.
			\end{itemize}
		\end{block}
		\vspace{0.02cm}
		\recallbox{\scriptsize \textbf{Key Rule:} Casino controls = \textbf{chip redemption + loyalty program tracking}.}
	\end{frame}
	
	% ================================================================
	% SLIDE 38: E-COMMERCE MONITORING — THE VOLUME TRAP
	% ================================================================
	\begin{frame}
		\frametitle{\fontsize{11}{13}\selectfont \textbf{E-Commerce Monitoring — The Volume Trap}}
		\begin{block}{\fontsize{7}{9}\selectfont \textbf{The Hard Question — CAMS Exam Favorite}}
			\scriptsize
			Why is transaction monitoring particularly challenging in an e-commerce context?
		\end{block}
		\vspace{0.02cm}
		\answerbox{\large \textbf{High transaction volumes and decentralized payments make patterns harder to detect.}}
		\vspace{0.02cm}
		\begin{block}{\fontsize{7}{9}\selectfont \textbf{Natural Explanation — Why You Got This Wrong}}
			\scriptsize
			\begin{itemize}
				\item \wronganswer{Transactions are always small and therefore low-risk} — \textbf{NO.} Small transactions can still be suspicious.
				\item \textbf{High volume:} Millions of transactions daily — traditional rules-based monitoring generates too many false positives.
				\item \textbf{Decentralized payments:} Multiple payment gateways, cryptocurrencies, digital wallets.
				\item \textbf{Key principle:} E-commerce challenge = \textbf{high volume + decentralized payments}.
			\end{itemize}
		\end{block}
		\vspace{0.02cm}
		\recallbox{\scriptsize \textbf{Key Rule:} E-commerce challenge = \textbf{high volume + decentralized payments}.}
	\end{frame}
	
	% ================================================================
	% SLIDE 39: DATA LINEAGE — THE TRANSPARENCY TRAP
	% ================================================================
	\begin{frame}
		\frametitle{\fontsize{11}{13}\selectfont \textbf{Data Lineage — The Transparency Trap}}
		\begin{block}{\fontsize{7}{9}\selectfont \textbf{The Hard Question — CAMS Exam Favorite}}
			\scriptsize
			Which support transparency and auditability through data lineage? (Select Two.)
		\end{block}
		\vspace{0.02cm}
		\answerbox{\large \textbf{1. Maintaining history of data transformations. 2. Capturing source of each data field.}}
		\vspace{0.02cm}
		\begin{block}{\fontsize{7}{9}\selectfont \textbf{Natural Explanation — Why You Got This Wrong}}
			\scriptsize
			\begin{itemize}
				\item \wronganswer{Archiving analyst notes from alert reviews} — This is \textbf{case documentation}, NOT data lineage.
				\item \wronganswer{Deleting outdated data annually} — \textbf{Undermines} auditability.
				\item \textbf{Data lineage:} Tracking the \textbf{flow of data} from source to destination.
				\item \textbf{History of transformations:} Shows how data was modified.
				\item \textbf{Source of each field:} Shows where data originated.
				\item \textbf{Key principle:} Data lineage = \textbf{transformations + source}. NOT case notes.
			\end{itemize}
		\end{block}
		\vspace{0.02cm}
		\recallbox{\scriptsize \textbf{Key Rule:} Data lineage = \textbf{transformations history + source of each field}.}
	\end{frame}
	
	% ================================================================
	% SLIDE 40: DATA VALIDATION — THE OBJECTIVES TRAP
	% ================================================================
	\begin{frame}
		\frametitle{\fontsize{11}{13}\selectfont \textbf{Data Validation — The Objectives Trap}}
		\begin{block}{\fontsize{7}{9}\selectfont \textbf{The Hard Question — CAMS Exam Favorite}}
			\scriptsize
			Which are typical objectives of data validation and testing in AFC systems? (Select Three.)
		\end{block}
		\vspace{0.02cm}
		\answerbox{\large \textbf{1. Ensuring data inputs reflect actual customer behavior. 2. Verifying automated controls function correctly. 3. Testing whether controls generate alerts under appropriate risk conditions.}}
		\vspace{0.02cm}
		\begin{block}{\fontsize{7}{9}\selectfont \textbf{Natural Explanation — Why You Got This Wrong}}
			\scriptsize
			\begin{itemize}
				\item \wronganswer{Confirming reports are generated in a format preferred by third-party vendors} — \textbf{NOT} a validation objective.
				\item \textbf{Data inputs reflect behavior:} Test that data from source systems accurately feeds into AML systems.
				\item \textbf{Controls functioning:} Validate that transaction monitoring rules fire alerts correctly.
				\item \textbf{Alerts under risk conditions:} Use known test cases to validate system logic.
				\item \textbf{Key principle:} Validation = \textbf{accuracy + functionality + appropriate alerts}.
			\end{itemize}
		\end{block}
		\vspace{0.02cm}
		\recallbox{\scriptsize \textbf{Key Rule:} Data validation = \textbf{accuracy + functionality + appropriate alerts}.}
	\end{frame}
	
	% ================================================================
	% SLIDE 41: CLUSTERING — THE BEHAVIORAL TRAP
	% ================================================================
	\begin{frame}
		\frametitle{\fontsize{11}{13}\selectfont \textbf{Clustering — The Behavioral Trap}}
		\begin{block}{\fontsize{7}{9}\selectfont \textbf{The Hard Question — CAMS Exam Favorite}}
			\scriptsize
			Which are characteristics of clustering in compliance analytics? (Select Two.)
		\end{block}
		\vspace{0.02cm}
		\answerbox{\large \textbf{1. It groups entities by shared behavioral characteristics. 2. It identifies connections across multiple transaction types.}}
		\vspace{0.02cm}
		\begin{block}{\fontsize{7}{9}\selectfont \textbf{Natural Explanation — Why You Got This Wrong}}
			\scriptsize
			\begin{itemize}
				\item \wronganswer{It excludes outliers to reduce financial crime risk} — \textbf{NO.} Outliers are often the \textbf{most suspicious}.
				\item \textbf{Clustering:} Groups entities that exhibit similar behaviors.
				\item \textbf{Purpose:} Detect networks of criminal activity (money mule networks, organized crime).
				\item \textbf{Connections across transaction types:} Links separate accounts that exhibit similar patterns.
				\item \textbf{Key principle:} Clustering = \textbf{shared behavior + connections}. NOT excluding outliers.
			\end{itemize}
		\end{block}
		\vspace{0.02cm}
		\recallbox{\scriptsize \textbf{Key Rule:} Clustering = \textbf{shared behaviors + connections}.}
	\end{frame}
	
	% ================================================================
	% SLIDE 42: DIGITAL ONBOARDING — THE BENEFITS TRAP
	% ================================================================
	\begin{frame}
		\frametitle{\fontsize{11}{13}\selectfont \textbf{Digital Onboarding — The Benefits Trap}}
		\begin{block}{\fontsize{7}{9}\selectfont \textbf{The Hard Question — CAMS Exam Favorite}}
			\scriptsize
			Which benefits are associated with digital onboarding technology according to FATF? (Select Two.)
		\end{block}
		\vspace{0.02cm}
		\answerbox{\large \textbf{1. Lower cost and increased auditability. 2. Greater accountability in compliance processes.}}
		\vspace{0.02cm}
		\begin{block}{\fontsize{7}{9}\selectfont \textbf{Natural Explanation — Why You Got This Wrong}}
			\scriptsize
			\begin{itemize}
				\item \wronganswer{Elimination of human review for high-risk customers} — \textbf{NO.} High-risk still require human review.
				\item \textbf{Lower cost:} Automation reduces manual effort.
				\item \textbf{Increased auditability:} Digital records are easier to track and audit.
				\item \textbf{Greater accountability:} Clear audit trails of who did what.
				\item \textbf{Key principle:} Digital onboarding = \textbf{lower cost + auditability + accountability}. NOT eliminating human review.
			\end{itemize}
		\end{block}
		\vspace{0.02cm}
		\recallbox{\scriptsize \textbf{Key Rule:} Digital onboarding = \textbf{lower cost + auditability + accountability}.}
	\end{frame}
	
	% ================================================================
	% SLIDE 43: PERPETUAL KYC — THE TRIGGER-BASED TRAP
	% ================================================================
	\begin{frame}
		\frametitle{\fontsize{11}{13}\selectfont \textbf{Perpetual KYC — The Trigger-Based Trap}}
		\begin{block}{\fontsize{7}{9}\selectfont \textbf{The Hard Question — CAMS Exam Favorite}}
			\scriptsize
			Which feature most effectively distinguishes perpetual KYC from traditional periodic reviews?
		\end{block}
		\vspace{0.02cm}
		\answerbox{\large \textbf{Trigger-based updates rather than fixed calendar reviews.}}
		\vspace{0.02cm}
		\begin{block}{\fontsize{7}{9}\selectfont \textbf{Natural Explanation — Why You Got This Wrong}}
			\scriptsize
			\begin{itemize}
				\item \wronganswer{It eliminates all manual review} — \textbf{NO.} Human review is still required.
				\item \textbf{Perpetual KYC (pKYC):} Continuous, event-driven KYC refresh.
				\item \textbf{Triggers:}
				\begin{itemize}
					\item Changes in PEP status
					\item Adverse media hits
					\item Unusual transaction patterns
					\item Jurisdictional risk changes
				\end{itemize}
				\item \textbf{Benefits:} Reduces customer friction, improves accuracy, ensures timely EDD.
				\item \textbf{Key principle:} pKYC = \textbf{trigger-based}, not calendar-based.
			\end{itemize}
		\end{block}
		\vspace{0.02cm}
		\recallbox{\scriptsize \textbf{Key Rule:} pKYC = \textbf{trigger-based updates}, not fixed calendar reviews.}
	\end{frame}
	
	% ================================================================
	% SLIDE 44: DEEPFAKE DETECTION — THE LIVENESS TRAP
	% ================================================================
	\begin{frame}
		\frametitle{\fontsize{11}{13}\selectfont \textbf{Deepfake Detection — The Liveness Trap}}
		\begin{block}{\fontsize{7}{9}\selectfont \textbf{The Scenario — CAMS Exam Favorite}}
			\scriptsize
			A fraudster used a high-resolution image and a deepfake video to bypass facial recognition.
		\end{block}
		\vspace{0.02cm}
		\begin{block}{\fontsize{7}{9}\selectfont \textbf{The Hard Question}}
			\scriptsize
			What tools could have prevented this? (Select Two.)
		\end{block}
		\vspace{0.02cm}
		\answerbox{\large \textbf{1. Software detecting inconsistencies in light or skin texture. 2. Prompts for the customer to blink, smile, or repeat a phrase.}}
		\vspace{0.02cm}
		\begin{block}{\fontsize{7}{9}\selectfont \textbf{Natural Explanation — Why You Got This Wrong}}
			\scriptsize
			\begin{itemize}
				\item \wronganswer{Real-time fraud risk scoring using behavioral analytics} — This is \textbf{behavioral analytics}, not deepfake detection.
				\item \textbf{Texture analysis:} Deepfakes often have inconsistencies in light, skin texture, and shadows.
				\item \textbf{Active liveness prompts:} Blink, smile, repeat a phrase — requires real-time action.
				\item \textbf{Key principle:} Deepfake detection = \textbf{texture analysis + active liveness prompts}.
			\end{itemize}
		\end{block}
		\vspace{0.02cm}
		\recallbox{\scriptsize \textbf{Key Rule:} Deepfake detection = \textbf{texture analysis + active liveness prompts}.}
	\end{frame}
	
	% ================================================================
	% SLIDE 45: AML COMMUNICATIONS — THE SCRIPT TRAP
	% ================================================================
	\begin{frame}
		\frametitle{\fontsize{11}{13}\selectfont \textbf{AML Communications — The Script Trap}}
		\begin{block}{\fontsize{7}{9}\selectfont \textbf{The Hard Question — CAMS Exam Favorite}}
			\scriptsize
			What should AML-sensitive communications emphasize to reduce risk? (Select Two.)
		\end{block}
		\vspace{0.02cm}
		\answerbox{\large \textbf{1. Avoid mentioning internal reviews. 2. Use pre-approved scripts for service disruptions.}}
		\vspace{0.02cm}
		\begin{block}{\fontsize{7}{9}\selectfont \textbf{Natural Explanation — Why You Got This Wrong}}
			\scriptsize
			\begin{itemize}
				\item \wronganswer{Refer customers to the AML department for financial advice} — \textbf{NO.} AML doesn't give financial advice.
				\item \textbf{Avoid mentioning internal reviews:} Prevents tipping off.
				\item \textbf{Pre-approved scripts:} Ensures consistent, compliant messaging.
				\item \textbf{Key principle:} Customer-facing staff must \textbf{not} reveal any AML investigations.
			\end{itemize}
		\end{block}
		\vspace{0.02cm}
		\recallbox{\scriptsize \textbf{Key Rule:} AML communications = \textbf{avoid reviews + use pre-approved scripts}.}
	\end{frame}
	
	% ================================================================
	% SLIDE 46: ENTITY RESOLUTION — THE BLOCKCHAIN TRAP
	% ================================================================
	\begin{frame}
		\frametitle{\fontsize{11}{13}\selectfont \textbf{Entity Resolution — The Blockchain Trap}}
		\begin{block}{\fontsize{7}{9}\selectfont \textbf{The Hard Question — CAMS Exam Favorite}}
			\scriptsize
			How does entity resolution improve blockchain analytics?
		\end{block}
		\vspace{0.02cm}
		\answerbox{\large \textbf{It links data from multiple sources to improve match accuracy.}}
		\vspace{0.02cm}
		\begin{block}{\fontsize{7}{9}\selectfont \textbf{Natural Explanation — Why You Got This Wrong}}
			\scriptsize
			\begin{itemize}
				\item \wronganswer{It replaces all manual reviews} — \textbf{NO.} It's a tool, not a replacement.
				\item \textbf{Entity resolution:} Links on-chain wallet addresses to off-chain risk data.
				\item \textbf{Multiple sources:} Wallet addresses, KYC data, sanctions lists, PEP lists, adverse media.
				\item \textbf{Improves match accuracy:} Reduces false positives and false negatives.
				\item \textbf{Key principle:} Entity resolution = \textbf{links data from multiple sources}.
			\end{itemize}
		\end{block}
		\vspace{0.02cm}
		\recallbox{\scriptsize \textbf{Key Rule:} Entity resolution = \textbf{links data from multiple sources} to improve accuracy.}
	\end{frame}
	
	% ================================================================
	% SLIDE 47: SAR/STR — THE THRESHOLD TRAP
	% ================================================================
	\begin{frame}
		\frametitle{\fontsize{11}{13}\selectfont \textbf{SAR/STR — The Threshold Trap}}
		\begin{block}{\fontsize{7}{9}\selectfont \textbf{The Hard Question — CAMS Exam Favorite}}
			\scriptsize
			Is there a minimum amount for SAR/STR filing?
		\end{block}
		\vspace{0.02cm}
		\answerbox{\large \textbf{No minimum amount — file based on reasonable suspicion.}}
		\vspace{0.02cm}
		\begin{block}{\fontsize{7}{9}\selectfont \textbf{Natural Explanation — Why You Got This Wrong}}
			\scriptsize
			\begin{itemize}
				\item \wronganswer{\$5,000 minimum} — That's a \textbf{US SAR threshold} (FIs), but NOT universal.
				\item \wronganswer{\$10,000 minimum} — That's \textbf{CTR}, not SAR.
				\item \textbf{SAR/STR = reasonable suspicion}, NOT amount.
				\item \textbf{Attempted transactions} must be reported.
				\item \textbf{Structuring} must be reported.
				\item \textbf{Key principle:} Suspicion drives filing, NOT amount.
			\end{itemize}
		\end{block}
		\vspace{0.02cm}
		\recallbox{\scriptsize \textbf{Key Rule:} SAR/STR = \textbf{reasonable suspicion}, NOT amount.}
	\end{frame}
	
	% ================================================================
	% SLIDE 48: SHELL COMPANIES — THE SECRECY LAWS TRAP
	% ================================================================
	\begin{frame}
		\frametitle{\fontsize{11}{13}\selectfont \textbf{Shell Companies — The Secrecy Laws Trap}}
		\begin{block}{\fontsize{7}{9}\selectfont \textbf{The Hard Question — CAMS Exam Favorite}}
			\scriptsize
			Which characteristics are commonly associated with shell companies used for ML? (Select Three.)
		\end{block}
		\vspace{0.02cm}
		\answerbox{\large \textbf{1. No significant assets or operations. 2. Registered in jurisdictions with strict secrecy laws. 3. Used to conceal beneficial ownership.}}
		\vspace{0.02cm}
		\begin{block}{\fontsize{7}{9}\selectfont \textbf{Natural Explanation — Why You Got This Wrong}}
			\scriptsize
			\begin{itemize}
				\item \wronganswer{Bearer shares issued to identifiable natural persons} — \textbf{NO.} Bearer shares are \textbf{anonymous}, not identifiable.
				\item \textbf{No assets/operations:} Shell companies are empty vehicles.
				\item \textbf{Strict secrecy laws:} Jurisdictions with secrecy laws enable anonymity.
				\item \textbf{Conceal BO:} The primary purpose is to hide who owns the company.
				\item \textbf{Key principle:} Shell companies = \textbf{empty + secretive + BO-concealing}.
			\end{itemize}
		\end{block}
		\vspace{0.02cm}
		\recallbox{\scriptsize \textbf{Key Rule:} Shell companies = \textbf{no operations + secrecy laws + BO concealment}.}
	\end{frame}
	
	% ================================================================
	% SLIDE 49: STRUCTURING — THE SMALL DEPOSITS TRAP
	% ================================================================
	\begin{frame}
		\frametitle{\fontsize{11}{13}\selectfont \textbf{Structuring — The Small Deposits Trap}}
		\begin{block}{\fontsize{7}{9}\selectfont \textbf{The Hard Question — CAMS Exam Favorite}}
			\scriptsize
			A fraud analyst discovers small cash deposits over several days, followed by a sudden international wire transfer. What laundering method is most likely being used?
		\end{block}
		\vspace{0.02cm}
		\answerbox{\large \textbf{Structuring (smurfing).}}
		\vspace{0.02cm}
		\begin{block}{\fontsize{7}{9}\selectfont \textbf{Natural Explanation — Why You Got This Wrong}}
			\scriptsize
			\begin{itemize}
				\item \wronganswer{Trade-based money laundering} — This involves \textbf{trade transactions}, not cash deposits.
				\item \wronganswer{Layering} — Layering involves \textbf{concealing origin}, but the initial pattern is structuring.
				\item \textbf{Structuring:} Breaking up large transactions into smaller amounts to avoid reporting thresholds.
				\item \textbf{Pattern:} Small deposits (avoid CTR threshold) + sudden wire transfer (integration).
				\item \textbf{Key principle:} Structuring = \textbf{small deposits to avoid thresholds} + sudden movement.
			\end{itemize}
		\end{block}
		\vspace{0.02cm}
		\recallbox{\scriptsize \textbf{Key Rule:} Structuring = \textbf{small deposits + sudden wire transfer}.}
	\end{frame}
	
	% ================================================================
	% SLIDE 50: TRADE-BASED ML — THE MISPRICING TRAP
	% ================================================================
	\begin{frame}
		\frametitle{\fontsize{11}{13}\selectfont \textbf{Trade-Based ML — The Mispricing Trap}}
		\begin{block}{\fontsize{7}{9}\selectfont \textbf{The Hard Question — CAMS Exam Favorite}}
			\scriptsize
			A textile exporter sends shipments priced significantly below market rates, and the counterpart pays the remainder via separate wire transfers to a third party. What does this suggest?
		\end{block}
		\vspace{0.02cm}
		\answerbox{\large \textbf{Intentional trade mispricing for money laundering.}}
		\vspace{0.02cm}
		\begin{block}{\fontsize{7}{9}\selectfont \textbf{Natural Explanation — Why You Got This Wrong}}
			\scriptsize
			\begin{itemize}
				\item \wronganswer{Normal business practice} — \textbf{NO.} Below-market pricing + third-party payments = suspicious.
				\item \wronganswer{Legitimate tax avoidance} — \textbf{NO.} Tax avoidance requires disclosure, not concealment.
				\item \textbf{Under-invoicing:} Goods priced below market to move value.
				\item \textbf{Third-party payments:} The difference is laundered through separate payments.
				\item \textbf{Key principle:} Trade mispricing = \textbf{under/over-invoicing + third-party payments}.
			\end{itemize}
		\end{block}
		\vspace{0.02cm}
		\recallbox{\scriptsize \textbf{Key Rule:} TBML mispricing = \textbf{below-market pricing + third-party payments}.}
	\end{frame}
	
	% ================================================================
	% SLIDE 51: DORA — THE CENTRALIZED OVERSIGHT TRAP
	% ================================================================
	\begin{frame}
		\frametitle{\fontsize{11}{13}\selectfont \textbf{DORA — The Centralized Oversight Trap}}
		\begin{block}{\fontsize{7}{9}\selectfont \textbf{The Hard Question — CAMS Exam Favorite}}
			\scriptsize
			A cross-border bank is reviewing digital risk governance. What action best aligns with DORA?
		\end{block}
		\vspace{0.02cm}
		\answerbox{\large \textbf{Centralizing oversight of ICT risk management across all member-state branches.}}
		\vspace{0.02cm}
		\begin{block}{\fontsize{7}{9}\selectfont \textbf{Natural Explanation — Why You Got This Wrong}}
			\scriptsize
			\begin{itemize}
				\item \wronganswer{Relying on each national regulator's cybersecurity requirements} — \textbf{Insufficient}. DORA mandates \textbf{pan-European} oversight.
				\item \textbf{Centralized} = unified governance across all EU operations.
				\item \textbf{Not} just compliance with each country's rules.
				\item Requires \textbf{single} ICT risk governance framework for the entire group.
				\item \textbf{Key principle:} DORA = \textbf{centralized ICT oversight}. National rules are not enough.
			\end{itemize}
		\end{block}
		\vspace{0.02cm}
		\recallbox{\scriptsize \textbf{Key Rule:} DORA = \textbf{centralized ICT oversight}.}
	\end{frame}
	
	% ================================================================
	% SLIDE 52: DEVICE INTELLIGENCE — THE SOCIAL MEDIA TRAP
	% ================================================================
	\begin{frame}
		\frametitle{\fontsize{11}{13}\selectfont \textbf{Device Intelligence — The Social Media Trap}}
		\begin{block}{\fontsize{7}{9}\selectfont \textbf{The Hard Question — CAMS Exam Favorite}}
			\scriptsize
			Which elements do device intelligence tools capture? (Select Three.)
		\end{block}
		\vspace{0.02cm}
		\answerbox{\large \textbf{1. IP address and geolocation. 2. Operating system and installed plugins. 3. Browser type and screen resolution.}}
		\vspace{0.02cm}
		\begin{block}{\fontsize{7}{9}\selectfont \textbf{Natural Explanation — Why You Got This Wrong}}
			\scriptsize
			\begin{itemize}
				\item \wronganswer{Social media accounts and browsing history} — \textbf{NO.} Device intelligence captures \textbf{device fingerprinting}, not social media.
				\item \textbf{IP address/geolocation:} Location and network information.
				\item \textbf{OS/plugins:} Device configuration.
				\item \textbf{Browser/screen resolution:} Browser fingerprinting.
				\item \textbf{Key principle:} Device intelligence = \textbf{device fingerprinting}, NOT social media or browsing history.
			\end{itemize}
		\end{block}
		\vspace{0.02cm}
		\recallbox{\scriptsize \textbf{Key Rule:} Device intelligence = \textbf{device fingerprinting}. NOT social media.}
	\end{frame}
	
	% ================================================================
	% SLIDE 53: AI TRANSITION — THE DATA/PEOPLE/EXPLAINABILITY TRAP
	% ================================================================
	\begin{frame}
		\frametitle{\fontsize{11}{13}\selectfont \textbf{AI Transition — The Data/People/Explainability Trap}}
		\begin{block}{\fontsize{7}{9}\selectfont \textbf{The Hard Question — CAMS Exam Favorite}}
			\scriptsize
			Which actions are most important when transitioning to AI-based compliance tools? (Select Three.)
		\end{block}
		\vspace{0.02cm}
		\answerbox{\large \textbf{1. Ensuring data quality for training algorithms. 2. Managing change in employee workflows. 3. Verifying transparency and explainability in AI models.}}
		\vspace{0.02cm}
		\begin{block}{\fontsize{7}{9}\selectfont \textbf{Natural Explanation — Why You Got This Wrong}}
			\scriptsize
			\begin{itemize}
				\item \wronganswer{Selecting third-party global audit providers} — This is \textbf{less critical} than data, people, and explainability.
				\item \textbf{Data quality:} "Garbage in, garbage out" — AI is only as good as its training data.
				\item \textbf{Employee workflows:} AI changes how analysts work — must provide adequate training.
				\item \textbf{Explainability:} Regulators require \textbf{explainable AI} — must justify why an AI flagged something.
				\item \textbf{Key principle:} AI transition = \textbf{data + people + explainability}.
			\end{itemize}
		\end{block}
		\vspace{0.02cm}
		\recallbox{\scriptsize \textbf{Key Rule:} AI transition = \textbf{data + people + explainability}.}
	\end{frame}
	
	% ================================================================
	% SLIDE 54: RPA — THE LARGE-VOLUME TRAP
	% ================================================================
	\begin{frame}
		\frametitle{\fontsize{11}{13}\selectfont \textbf{RPA — The Large-Volume Trap}}
		\begin{block}{\fontsize{7}{9}\selectfont \textbf{The Hard Question — CAMS Exam Favorite}}
			\scriptsize
			One advantage of RPA in AML operations is that it can:
		\end{block}
		\vspace{0.02cm}
		\answerbox{\large \textbf{Reduce compliance costs by handling large-volume, repeatable tasks.}}
		\vspace{0.02cm}
		\begin{block}{\fontsize{7}{9}\selectfont \textbf{Natural Explanation — Why You Got This Wrong}}
			\scriptsize
			\begin{itemize}
				\item \wronganswer{Replace all human decision-making} — \textbf{NO.} RPA automates \textbf{tasks}, not decisions.
				\item \wronganswer{Eliminate all false positives} — \textbf{NO.} RPA doesn't eliminate false positives.
				\item \textbf{RPA:} Handles large-volume, repeatable tasks (data extraction, reconciliation, report generation).
				\item \textbf{Key principle:} RPA = \textbf{large-volume, repeatable tasks}. NOT decision-making.
			\end{itemize}
		\end{block}
		\vspace{0.02cm}
		\recallbox{\scriptsize \textbf{Key Rule:} RPA = \textbf{large-volume, repeatable tasks}.}
	\end{frame}
	
	% ================================================================
	% SLIDE 55: BIOMETRICS — THE COMPROMISED DATA TRAP
	% ================================================================
	\begin{frame}
		\frametitle{\fontsize{11}{13}\selectfont \textbf{Biometrics — The Compromised Data Trap}}
		\begin{block}{\fontsize{7}{9}\selectfont \textbf{The Hard Question — CAMS Exam Favorite}}
			\scriptsize
			Which risk is specifically associated with facial recognition compared to passwords?
		\end{block}
		\vspace{0.02cm}
		\answerbox{\large \textbf{Compromised biometric data is difficult to change or reset unlike passwords.}}
		\vspace{0.02cm}
		\begin{block}{\fontsize{7}{9}\selectfont \textbf{Natural Explanation — Why You Got This Wrong}}
			\scriptsize
			\begin{itemize}
				\item \wronganswer{The technology requires more storage capacity} — Not a unique risk.
				\item \wronganswer{Facial recognition is less secure} — Not necessarily true.
				\item \textbf{Passwords can be changed} — if compromised, you reset them.
				\item \textbf{Biometrics cannot be changed} — you can't get a new face or fingerprint.
				\item Once biometric data is compromised, it's \textbf{permanently compromised}.
				\item \textbf{Key principle:} Biometric risk = \textbf{irreplaceable} — can't change like passwords.
			\end{itemize}
		\end{block}
		\vspace{0.02cm}
		\recallbox{\scriptsize \textbf{Key Rule:} Biometric risk = \textbf{can't change like passwords}.}
	\end{frame}
	
	% ================================================================
	% SLIDE 56: NATIONAL e-ID — THE INTEGRATION TRAP
	% ================================================================
	\begin{frame}
		\frametitle{\fontsize{11}{13}\selectfont \textbf{National e-ID — The Integration Trap}}
		\begin{block}{\fontsize{7}{9}\selectfont \textbf{The Hard Question — CAMS Exam Favorite}}
			\scriptsize
			A fintech relies solely on device recognition for authentication. What additional measures could improve controls? (Select Two.)
		\end{block}
		\vspace{0.02cm}
		\answerbox{\large \textbf{1. Implementing document verification with biometric authentication. 2. Integrating with national electronic identity verification systems.}}
		\vspace{0.02cm}
		\begin{block}{\fontsize{7}{9}\selectfont \textbf{Natural Explanation — Why You Got This Wrong}}
			\scriptsize
			\begin{itemize}
				\item \wronganswer{Enhancing digital footprint analysis through IP verification} — IP verification is \textbf{already part} of device recognition.
				\item \wronganswer{Relying solely on social media verification} — Not a reliable identity verification method.
				\item \textbf{Document verification + biometrics:} Ensures identity authenticity.
				\item \textbf{National e-ID integration:} Leverages trusted government identity databases.
				\item \textbf{Key principle:} Identity verification = \textbf{document verification + biometrics + e-ID integration}.
			\end{itemize}
		\end{block}
		\vspace{0.02cm}
		\recallbox{\scriptsize \textbf{Key Rule:} Identity verification = \textbf{document + biometrics + e-ID}.}
	\end{frame}
	
	% ================================================================
	% SLIDE 57: OFAC SANCTIONS SCREENING — THE PROPORTIONALITY TRAP
	% ================================================================
	\begin{frame}
		\frametitle{\fontsize{11}{13}\selectfont \textbf{OFAC Screening — The Proportionality Trap}}
		\begin{block}{\fontsize{7}{9}\selectfont \textbf{The Hard Question — CAMS Exam Favorite}}
			\scriptsize
			A luxury brand with modest client base wants sanctions screening. Which setup is most appropriate?
		\end{block}
		\vspace{0.02cm}
		\answerbox{\large \textbf{Using spreadsheets linked to the OFAC website for customer checks.}}
		\vspace{0.02cm}
		\begin{block}{\fontsize{7}{9}\selectfont \textbf{Natural Explanation — Why You Got This Wrong}}
			\scriptsize
			\begin{itemize}
				\item \wronganswer{Implementing a batch screening audit conducted annually} — This is \textbf{reactive}, not proactive.
				\item \wronganswer{Full automated real-time system} — \textbf{Overkill} for a small business; cost-prohibitive.
				\item \textbf{Spreadsheets + OFAC website:} Proportionate for small volume.
				\item \textbf{Risk-based approach:} Allows scaled controls based on size and risk.
				\item \textbf{Key principle:} Screening must be \textbf{proportionate} to risk and size.
			\end{itemize}
		\end{block}
		\vspace{0.02cm}
		\recallbox{\scriptsize \textbf{Key Rule:} Screening = \textbf{proportionate to risk and size}.}
	\end{frame}
	
	% ================================================================
	% SLIDE 58: RULES-BASED MONITORING — THE LEGACY CUSTOMER TRAP
	% ================================================================
	\begin{frame}
		\frametitle{\fontsize{11}{13}\selectfont \textbf{Rules-Based Monitoring — The Legacy Customer Trap}}
		\begin{block}{\fontsize{7}{9}\selectfont \textbf{The Hard Question — CAMS Exam Favorite}}
			\scriptsize
			A rules-based system generates high alerts for a legacy customer whose monthly deposits have been consistent for years. What is the most appropriate next step?
		\end{block}
		\vspace{0.02cm}
		\answerbox{\large \textbf{Update customer information about expected transaction patterns in the monitoring system.}}
		\vspace{0.02cm}
		\begin{block}{\fontsize{7}{9}\selectfont \textbf{Natural Explanation — Why You Got This Wrong}}
			\scriptsize
			\begin{itemize}
				\item \wronganswer{Escalate all alerts to Level 2} — This is \textbf{inefficient} and wastes resources.
				\item \wronganswer{Increase the alert threshold} — Could \textbf{miss} real risk.
				\item \textbf{Update expected patterns:} The system needs to know what's \textbf{normal} for this customer.
				\item This reduces false positives \textbf{while maintaining detection}.
				\item \textbf{Key principle:} Monitoring systems need \textbf{up-to-date customer profiles}.
			\end{itemize}
		\end{block}
		\vspace{0.02cm}
		\recallbox{\scriptsize \textbf{Key Rule:} Update expected patterns to reduce false positives.}
	\end{frame}
	
	% ================================================================
	% SLIDE 59: G-20 ANTI-CORRUPTION — THE ACTION PLAN TRAP
	% ================================================================
	\begin{frame}
		\frametitle{\fontsize{11}{13}\selectfont \textbf{G-20 Anti-Corruption — The Action Plan Trap}}
		\begin{block}{\fontsize{7}{9}\selectfont \textbf{The Hard Question — CAMS Exam Favorite}}
			\scriptsize
			The G-20 Anti-Corruption Working Group contributes to anti-corruption efforts by:
		\end{block}
		\vspace{0.02cm}
		\answerbox{\large \textbf{Coordinating with member states to establish and implement multi-year anti-corruption action plans.}}
		\vspace{0.02cm}
		\begin{block}{\fontsize{7}{9}\selectfont \textbf{Natural Explanation — Why You Got This Wrong}}
			\scriptsize
			\begin{itemize}
				\item \wronganswer{Setting CPI scores based on national prosecutions} — That's \textbf{Transparency International}, not G-20.
				\item \textbf{G-20 ACWG} = coordination and action plans.
				\item \textbf{Multi-year action plans:} Structured, long-term anti-corruption efforts.
				\item \textbf{Key principle:} G-20 ACWG = \textbf{coordination + action plans}. NOT CPI.
			\end{itemize}
		\end{block}
		\vspace{0.02cm}
		\recallbox{\scriptsize \textbf{Key Rule:} G-20 ACWG = \textbf{coordination + action plans}.}
	\end{frame}
	
	% ================================================================
	% SLIDE 60: BASEL AML INDEX — THE VULNERABILITY TRAP
	% ================================================================
	\begin{frame}
		\frametitle{\fontsize{11}{13}\selectfont \textbf{Basel AML Index — The Vulnerability Trap}}
		\begin{block}{\fontsize{7}{9}\selectfont \textbf{The Hard Question — CAMS Exam Favorite}}
			\scriptsize
			Why would a national enforcement agency choose the Basel AML Index?
		\end{block}
		\vspace{0.02cm}
		\answerbox{\large \textbf{It compiles vulnerability indicators across legal, political, and financial dimensions.}}
		\vspace{0.02cm}
		\begin{block}{\fontsize{7}{9}\selectfont \textbf{Natural Explanation — Why You Got This Wrong}}
			\scriptsize
			\begin{itemize}
				\item \wronganswer{Ranks countries by GDP growth rate} — \textbf{NO.} That's economic, not AML.
				\item \wronganswer{Tracks individual ML convictions} — Too narrow; not the purpose.
				\item \textbf{Basel AML Index} = \textbf{country risk assessment} for AML/CFT.
				\item Compiles \textbf{vulnerability indicators} across multiple dimensions.
				\item Used for \textbf{due diligence} and \textbf{country risk scoring}.
				\item \textbf{Key principle:} Basel AML Index = \textbf{vulnerability indicators}.
			\end{itemize}
		\end{block}
		\vspace{0.02cm}
		\recallbox{\scriptsize \textbf{Key Rule:} Basel AML Index = \textbf{vulnerability indicators}.}
	\end{frame}
	
	% ================================================================
	% SLIDE 61: FINAL EXAM TIPS — ALL CONCEPTS SUMMARY
	% ================================================================
	\begin{frame}
		\frametitle{\fontsize{11}{13}\selectfont \textbf{FINAL EXAM TIPS — All Concepts Summary}}
		\scriptsize
		\begin{block}{\fontsize{7}{9}\selectfont \textbf{Top Traps — Natural Language Summary}}
			\scriptsize
			\begin{enumerate}
				\item \textbf{Authentication:} "Password biometrics" = trap. Use MFA + physiological + behavioral.
				\item \textbf{VASP/MiCA:} White paper + one notification. NOT each country.
				\item \textbf{MLAT:} Challenge = admissibility, not just sharing.
				\item \textbf{PPP:} Real-time two-way communication. NOT post-hoc reporting.
				\item \textbf{FATF 6-Step:} Environmental + horizon scanning + threat identification.
				\item \textbf{FSI:} BO reporting + cross-border services + secrecy-to-scale ratio.
				\item \textbf{EWRA:} MLRO shapes risk. Internal Audit assesses controls.
				\item \textbf{Tax:} Avoidance = legal with disclosure. Evasion = illegal with concealment.
				\item \textbf{Board:} Oversight and challenge. NOT operational involvement.
				\item \textbf{TF:} Fake charities + Hawala + physical cash. NOT regulated wire transfers.
				\item \textbf{Trafficking:} Mobile payments + weak cross-border enforcement.
				\item \textbf{TCSP:} Cannot outsource AML responsibility.
				\item \textbf{DNFBP:} Committee + questionnaires. NOT blanket EDD.
				\item \textbf{Layering:} Concealing origin. Placement = entering system.
				\item \textbf{Proactive:} Update EWRA immediately.
				\item \textbf{IOSCO:} Investor protection + financial stability.
				\item \textbf{Open Source:} Entity-specific research. NOT general guidance.
				\item \textbf{Batch Screening:} Large volumes + changing risks. NOT real-time.
				\item \textbf{Data Extraction:} Normalization challenge.
				\item \textbf{Data Mining:} Internal + external sources + intelligent filters.
				\item \textbf{Data Integration:} Connection, not isolation.
				\item \textbf{BSA:} Tailoring + assigning roles. NOT fixed thresholds.
				\item \textbf{China AML 2025:} Expanded to DNFBPs.
				\item \textbf{AMLA:} Cross-border, high-risk = AMLA oversight.
				\item \textbf{DORA:} Centralized ICT oversight.
				\item \textbf{AI Transition:} Data + people + explainability.
				\item \textbf{RPA:} Large-volume repeatable tasks.
				\item \textbf{Entity Resolution:} Links data from multiple sources.
				\item \textbf{SAR/STR:} Reasonable suspicion, NOT amount.
				\item \textbf{Structuring:} Small deposits + sudden wire transfer.
				\item \textbf{TBML:} Below-market pricing + third-party payments.
			\end{enumerate}
		\end{block}
		\vspace{0.02cm}
		\recallbox{\scriptsize \textbf{Final Rule:} Understand the \textbf{why} behind each concept. The exam tests \textbf{application}, not memorization.}
	\end{frame}
	
	% ================================================================
	% END SLIDE
	% ================================================================
	\begin{frame}
		\centering
		\vspace{1.5cm}
		{\fontsize{18}{22}\selectfont \textbf{CAMS EXAM — COMPLETE}}\\[0.2cm]
		{\fontsize{11}{13}\selectfont All Hard Questions \& Exam Traps Explained}\\[0.1cm]
		\rule{4cm}{0.5pt}\\[0.1cm]
		{\fontsize{8}{10}\selectfont 60+ Concepts · All Domains · Natural Explanations}\\[0.1cm]
		\rule{4cm}{0.5pt}\\[0.1cm]
		{\fontsize{8}{10}\selectfont \textit{Ready for the exam. Good luck!}}
		\vspace{0.5cm}
	\end{frame}
	
\end{document}